\documentclass[8pt]{beamer}
\usetheme{Madrid}
\usecolortheme{default}
\usepackage{amsmath, amssymb, graphicx, array, multirow, booktabs, xcolor, tikz}
\usetikzlibrary{shapes,arrows,positioning,fit,backgrounds}

% Compact font settings
\setbeamerfont{title}{size=\Large}
\setbeamerfont{subtitle}{size=\small}
\setbeamerfont{author}{size=\scriptsize}
\setbeamerfont{date}{size=\scriptsize}
\setbeamerfont{frametitle}{size=\large}
\setbeamerfont{framesubtitle}{size=\normalsize}
\setbeamerfont{enumerate item}{size=\footnotesize}
\setbeamerfont{itemize item}{size=\footnotesize}
\setbeamerfont{block title}{size=\footnotesize}

% Compact spacing
\setlength{\itemsep}{0.08ex}
\setlength{\parskip}{0.08ex}
\setlength{\leftmargini}{0.3cm}
\setbeamersize{text margin left=3mm, text margin right=3mm}
\addtobeamertemplate{frame begin}{\vspace{-0.4cm}}{}

% Colors
\definecolor{myblue}{RGB}{0,0,255}
\definecolor{myred}{RGB}{255,0,0}
\definecolor{mygreen}{RGB}{0,128,0}
\definecolor{myorange}{RGB}{255,165,0}
\definecolor{mypurple}{RGB}{128,0,128}
\definecolor{mygray}{RGB}{128,128,128}
\definecolor{mygold}{RGB}{255,215,0}

\title{Module 3: Risk and Risk Factors}
\subtitle{100 Tutorial Topics: Algorithmic Bias, Fairness, Explainability, GenAI Risks}
\author{GARP RAI Exam Preparation}
\date{}

\begin{document}
	
	\begin{frame}
		\titlepage
	\end{frame}
	
	% ============================================================
	% SECTION 1: INTRODUCTION TO AI RISKS
	% ============================================================
	
	\begin{frame}{T1: Introduction to AI Risks}
		\fontsize{6.5}{7.5}\selectfont
		\textbf{TOPIC: Why AI Risk Management is Challenging}
		
		\vspace{0.1cm}
		\textbf{KEY CHARACTERISTICS OF AI RISKS:}
		\begin{itemize}
			\item \textcolor{myblue}{Pervasiveness:} AI is used across many domains - risk takes different forms depending on use case
			\item \textcolor{myblue}{High-Stakes Capabilities:} AI increasingly used in consequential settings - shortcomings have significant consequences
			\item \textcolor{myblue}{Rapid Development:} Capabilities advance faster than understanding of risks
			\item \textcolor{myblue}{Governance Challenges:} Regulatory responses are slow to develop and difficult to change
		\end{itemize}
		
		\vspace{0.1cm}
		\textbf{WHY IT MATTERS:}
		\begin{itemize}
			\item Loan application algorithms profoundly affect applicants
			\item Medical diagnostic algorithms have long-lasting impact on patients
			\item Slow regulation allows harmful technologies to persist
			\item Normalization and entrenchment of harmful practices
		\end{itemize}
		
		\vspace{0.1cm}
		\textbf{INTERCONNECTED NATURE OF RISKS:}
		\begin{itemize}
			\item Lack of transparency affects fairness and safety
			\item Algorithmic bias leads to reputational risks
			\item Risks reinforce each other - need holistic approach
		\end{itemize}
		
		\textcolor{myred}{\textbf{EXAM TRAP:}} AI risks are \textcolor{myblue}{evolving, pervasive, and interconnected} - not static or isolated!
	\end{frame}
	
	\begin{frame}{T2: Pervasiveness Challenge}
		\fontsize{6.5}{7.5}\selectfont
		\textbf{TOPIC: Pervasiveness of AI and Risk Prediction Challenges}
		
		\vspace{0.1cm}
		\textbf{DEFINITION:}
		\begin{itemize}
			\item AI is hugely pervasive in both range and variety of use
			\item Applications span many domains: finance, healthcare, justice, employment, etc.
		\end{itemize}
		
		\vspace{0.1cm}
		\textbf{WHY PERVASIVENESS POSES CHALLENGES:}
		\begin{itemize}
			\item Risks can take on different forms depending on the use case
			\item A loan denial algorithm has different risks than a medical diagnostic algorithm
			\item Risk prediction must account for context
			\item No one-size-fits-all approach
		\end{itemize}
		
		\vspace{0.1cm}
		\textbf{EXAMPLES OF DIVERSE APPLICATIONS:}
		\begin{itemize}
			\item \textcolor{myblue}{Credit Scoring:} Assessing loan applications
			\item \textcolor{myblue}{Medical Diagnostics:} Identifying diseases from images
			\item \textcolor{myblue}{Policing:} Predictive policing and bail decisions
			\item \textcolor{myblue}{Employment:} Resume screening and hiring
			\item \textcolor{myblue}{Education:} College admissions
		\end{itemize}
		
		\textcolor{myred}{\textbf{EXAM TRAP:}} AI risks are \textcolor{myblue}{context-dependent} - must understand specific use case!
	\end{frame}
	
	\begin{frame}{T3: High-Stakes AI Use}
		\fontsize{6.5}{7.5}\selectfont
		\textbf{TOPIC: Consequences of AI in High-Stakes Settings}
		
		\vspace{0.1cm}
		\textbf{DEFINITION:}
		\begin{itemize}
			\item AI technologies have increasingly large capabilities
			\item Used in increasingly high-stakes settings
			\item Technological shortcomings can have significant consequences for people
		\end{itemize}
		
		\vspace{0.1cm}
		\textbf{EXAMPLES OF HIGH-STAKES CONSEQUENCES:}
		\begin{itemize}
			\item \textcolor{myblue}{Loan Applications:} Profound effect on applicant's financial well-being
			\item \textcolor{myblue}{Medical Diagnosis:} Long-lasting impact on patient's life
			\item \textcolor{myblue}{Bail Decisions:} Affect personal freedom
			\item \textcolor{myblue}{Hiring Decisions:} Impact career trajectories
		\end{itemize}
		
		\vspace{0.1cm}
		\textbf{IMPLICATIONS:}
		\begin{itemize}
			\item Algorithmic errors can have profound impacts on people's lives
			\item Need for robust testing and validation
			\item Human oversight essential
		\end{itemize}
		
		\textcolor{myred}{\textbf{EXAM TRAP:}} High-stakes AI = \textcolor{myblue}{significant consequences for individuals}!
	\end{frame}
	
	\begin{frame}{T4: Rapid AI Development}
		\fontsize{6.5}{7.5}\selectfont
		\textbf{TOPIC: Challenges from Rapid AI Development}
		
		\vspace{0.1cm}
		\textbf{DEFINITION:}
		\begin{itemize}
			\item AI development is rapid
			\item Capabilities advance faster than our understanding of risks
		\end{itemize}
		
		\vspace{0.1cm}
		\textbf{CHALLENGES CREATED:}
		\begin{itemize}
			\item \textcolor{myblue}{Impact Prediction:} Hard to predict future capabilities and risks
			\item \textcolor{myblue}{Risk Measurement:} Metrics may not capture new risks
			\item \textcolor{myblue}{Risk Mitigation:} Solutions may lag behind new developments
			\item \textcolor{myblue}{Governance:} Difficult to establish robust mechanisms
		\end{itemize}
		
		\vspace{0.1cm}
		\textbf{REGULATORY CHALLENGES:}
		\begin{itemize}
			\item Regulatory responses are slow to develop
			\item Difficult to change once established
			\item Slow regulation allows harmful technologies to persist
			\item Normalization and entrenchment of harmful practices
		\end{itemize}
		
		\textcolor{myred}{\textbf{EXAM TRAP:}} Rapid development = \textcolor{myblue}{governance and regulation challenges}!
	\end{frame}
	
	\begin{frame}{T5: Interconnected AI Risks}
		\fontsize{6.5}{7.5}\selectfont
		\textbf{TOPIC: How AI Risks are Interconnected}
		
		\vspace{0.1cm}
		\textbf{DEFINITION:}
		\begin{itemize}
			\item Risks from AI are often interconnected
			\item Addressing one risk may affect others
		\end{itemize}
		
		\vspace{0.1cm}
		\textbf{EXAMPLES OF INTERCONNECTIONS:}
		\begin{itemize}
			\item \textcolor{myblue}{Transparency $\to$ Fairness:} Lack of transparency affects fairness assessment
			\item \textcolor{myblue}{Bias $\to$ Reputation:} Algorithmic bias leads to reputational risks
			\item \textcolor{myblue}{Safety $\to$ Fairness:} Safety failures can reveal fairness issues
			\item \textcolor{myblue}{Explainability $\to$ Accountability:} Lack of explainability creates accountability gaps
		\end{itemize}
		
		\vspace{0.1cm}
		\textbf{IMPLICATIONS:}
		\begin{itemize}
			\item Need holistic approach to risk management
			\item Cannot address risks in isolation
			\item Must consider entire AI lifecycle
		\end{itemize}
		
		\textcolor{myred}{\textbf{EXAM TRAP:}} AI risks are \textcolor{myblue}{interconnected and mutually reinforcing}!
	\end{frame}
	
	% ============================================================
	% SECTION 2: ALGORITHMIC BIAS AND FAIRNESS
	% ============================================================
	
	\begin{frame}{T6: Algorithmic Bias Definition}
		\fontsize{6.5}{7.5}\selectfont
		\textbf{TOPIC: Algorithmic Bias - Definition and Types}
		
		\vspace{0.1cm}
		\textbf{DEFINITION:}
		\begin{itemize}
			\item Algorithmic bias is a systematic deviation of an algorithm's output, performance, or impact relative to some norm, aim, standard, or baseline
			\item Systematic (not random)
			\item Can affect output, performance, or impact
			\item Relative to some norm, aim, standard, or baseline
		\end{itemize}
		
		\vspace{0.1cm}
		\textbf{TWO TYPES:}
		\begin{itemize}
			\item \textcolor{myblue}{Explicit Bias:} Developer intentionally projects own biases into algorithm
			\item \textcolor{myblue}{Implicit Bias:} Bias creeps in undetected through design or data
		\end{itemize}
		
		\vspace{0.1cm}
		\textbf{WHY IT MATTERS:}
		\begin{itemize}
			\item Algorithms can exhibit bias and create unfair results
			\item Systematic nature means bias affects many people consistently
			\item Regulators increasingly demand fair and non-discriminatory AI
		\end{itemize}
		
		\textcolor{myred}{\textbf{EXAM TRAP:}} Algorithmic bias = \textcolor{myblue}{systematic deviation from a norm}!
	\end{frame}
	
	\begin{frame}{T7: Explicit vs Implicit Bias}
		\fontsize{6.5}{7.5}\selectfont
		\textbf{TOPIC: Distinguishing Between Explicit and Implicit Bias}
		
		\vspace{0.1cm}
		\textbf{EXPLICIT BIAS:}
		\begin{itemize}
			\item Developer intentionally projects own biases into algorithm
			\item Leads to biased outputs
			\item Intentional discrimination
		\end{itemize}
		
		\vspace{0.1cm}
		\textbf{IMPLICIT BIAS:}
		\begin{itemize}
			\item Bias creeps in undetected
			\item Through algorithm's design or training data
			\item Unintended
			\item Often harder to detect and correct
		\end{itemize}
		
		\vspace{0.1cm}
		\textbf{EXAMPLES:}
		\begin{itemize}
			\item \textcolor{myblue}{Explicit:} Deliberately excluding certain groups
			\item \textcolor{myblue}{Implicit:} Historical data reflecting past discrimination
		\end{itemize}
		
		\vspace{0.1cm}
		\textbf{KEY INSIGHT:}
		\begin{itemize}
			\item Both types can lead to unfair outcomes
			\item Implicit bias is more common in practice
			\item Need to check for both types
		\end{itemize}
		
		\textcolor{myred}{\textbf{EXAM TRAP:}} Explicit = \textcolor{myblue}{intentional}; Implicit = \textcolor{myblue}{unintended}!
	\end{frame}
	
	\begin{frame}{T8: Individual Fairness}
		\fontsize{6.5}{7.5}\selectfont
		\textbf{TOPIC: Individual Fairness - Treating Like Cases Alike}
		
		\vspace{0.1cm}
		\textbf{DEFINITION:}
		\begin{itemize}
			\item Based on Aristotelian doctrine of "treating like cases alike"
			\item Similarly situated individuals should be treated equally
			\item Focuses on individuals, not groups
		\end{itemize}
		
		\vspace{0.1cm}
		\textbf{KEY CHARACTERISTICS:}
		\begin{itemize}
			\item Treating applicants alike in all relevant aspects
			\item Avoids direct discrimination on protected characteristics
			\item Requires agreement on what counts as "alike"
		\end{itemize}
		
		\vspace{0.1cm}
		\textbf{LIMITATIONS:}
		\begin{itemize}
			\item What counts as "alike" is often contentious
			\item Should we consider hardships overcome?
			\item How to translate resilience into numbers?
			\item Different people have different views
		\end{itemize}
		
		\vspace{0.1cm}
		\textbf{EXAMPLE (College Admission):}
		\begin{itemize}
			\item If grades are the only relevant metric, select highest grades
			\item But students start from different positions
			\item Should socioeconomic factors matter?
		\end{itemize}
		
		\textcolor{myred}{\textbf{EXAM TRAP:}} Individual fairness = \textcolor{myblue}{treating like cases alike}!
	\end{frame}
	
	\begin{frame}{T9: Group Fairness}
		\fontsize{6.5}{7.5}\selectfont
		\textbf{TOPIC: Group Fairness - Statistical Differences Across Groups}
		
		\vspace{0.1cm}
		\textbf{DEFINITION:}
		\begin{itemize}
			\item Does not consider individuals
			\item Looks at statistical differences across groups
			\item Groups share protected characteristics
		\end{itemize}
		
		\vspace{0.1cm}
		\textbf{KEY CHARACTERISTICS:}
		\begin{itemize}
			\item Examines whether groups are disadvantaged
			\item Examples: admission rates by gender, loan approvals by race
			\item Two categories:
			\begin{itemize}
				\item Performance equality
				\item Output distribution
			\end{itemize}
		\end{itemize}
		
		\vspace{0.1cm}
		\textbf{MAIN MEASURES:}
		\begin{itemize}
			\item \textcolor{myblue}{Demographic Parity:} Equal distribution of predictions
			\item \textcolor{myblue}{Predictive Rate Parity:} Equal precision
			\item \textcolor{myblue}{Equal Opportunity:} Equal sensitivity
		\end{itemize}
		
		\textcolor{myred}{\textbf{EXAM TRAP:}} Group fairness = \textcolor{myblue}{statistical differences across groups}!
	\end{frame}
	
	\begin{frame}{T10: College Admission Example}
		\fontsize{6.5}{7.5}\selectfont
		\textbf{TOPIC: Fairness Challenges in College Admission}
		
		\vspace{0.1cm}
		\textbf{SCENARIO:}
		\begin{itemize}
			\item College uses algorithm CollegeAd
			\item Analyzes test scores, transcripts, writing samples
			\item Outputs hierarchy of most promising students
		\end{itemize}
		
		\vspace{0.1cm}
		\textbf{FAIRNESS QUESTIONS:}
		\begin{itemize}
			\item What counts as "alike"?
			\item Should we consider hardships overcome?
			\item How to translate resilience into numbers?
			\item Should socioeconomic factors matter?
		\end{itemize}
		
		\vspace{0.1cm}
		\textbf{LIMITATIONS OF INDIVIDUAL FAIRNESS:}
		\begin{itemize}
			\item Intuitive but requires agreement on "alike"
			\item Often contentious
			\item Designers must be explicit about relevant factors
			\item Factors may vary with context
		\end{itemize}
		
		\textcolor{myred}{\textbf{EXAM TRAP:}} Individual fairness requires \textcolor{myblue}{agreement on what is "alike"}!
	\end{frame}
	
	\begin{frame}{T11: Demographic Parity}
		\fontsize{6.5}{7.5}\selectfont
		\textbf{TOPIC: Demographic Parity - Definition and Formula}
		
		\vspace{0.1cm}
		\textbf{DEFINITION:}
		\begin{itemize}
			\item Distribution of predictions is identical across subpopulations
			\item Admission rate same across groups
			\item Loan approval rate same across groups
		\end{itemize}
		
		\vspace{0.1cm}
		\textbf{FORMULA:}
		\[
		P(\hat{Y}=1|A=0) = P(\hat{Y}=1|A=1)
		\]
		where:
		\begin{itemize}
			\item $\hat{Y} \in \{0,1\}$: Predicted outcome
			\item $A \in \{0,1\}$: Protected attribute/group membership
		\end{itemize}
		
		\vspace{0.1cm}
		\textbf{KEY CHARACTERISTIC:}
		\begin{itemize}
			\item Entirely independent of true value
			\item Does not consider if predictions are correct
			\item Only if predictions are equally distributed
		\end{itemize}
		
		\textcolor{myred}{\textbf{EXAM TRAP:}} Demographic parity = \textcolor{myblue}{equal distribution of predictions}!
	\end{frame}
	
	\begin{frame}{T12: Demographic Parity - Drawbacks}
		\fontsize{6.5}{7.5}\selectfont
		\textbf{TOPIC: Drawbacks of Demographic Parity}
		
		\vspace{0.1cm}
		\textbf{MAJOR DRAWBACK:}
		\begin{itemize}
			\item Does not consider quality/accuracy of predictions
			\item Can be satisfied even if all predictions are wrong
			\item Example: Algorithm performs worse on subpopulation, labels many unqualified as qualified
			\item Still satisfies demographic parity if proportion is same
		\end{itemize}
		
		\vspace{0.1cm}
		\textbf{WHEN INAPPROPRIATE:}
		\begin{itemize}
			\item Base rates significantly differ between groups
			\item Differences are relevant to decision
			\item Medical context: quality of prediction is primary concern
			\item Credit risk: relevant factors correlated with protected attributes
		\end{itemize}
		
		\vspace{0.1cm}
		\textbf{WHEN USEFUL:}
		\begin{itemize}
			\item When true value not directly accessible
			\item Credit scoring: time lapse before outcome observed
			\item Fraud detection: true fraud cases limited
		\end{itemize}
		
		\textcolor{myred}{\textbf{EXAM TRAP:}} Demographic parity ignores \textcolor{myblue}{prediction accuracy}!
	\end{frame}
	
	\begin{frame}{T13: Predictive Rate Parity}
		\fontsize{6.5}{7.5}\selectfont
		\textbf{TOPIC: Predictive Rate Parity - Definition and Formula}
		
		\vspace{0.1cm}
		\textbf{DEFINITION:}
		\begin{itemize}
			\item Precision is equal across subpopulations
			\item Among those predicted positive, proportion correctly identified is same
			\item Also called Positive Predictive Value (PPV) parity
		\end{itemize}
		
		\vspace{0.1cm}
		\textbf{FORMULA:}
		\[
		P(Y=1|\hat{Y}=1,A=0) = P(Y=1|\hat{Y}=1,A=1)
		\]
		
		\vspace{0.1cm}
		\textbf{IN TERMS OF PRECISION:}
		\[
		\text{Precision}_0 = \text{Precision}_1
		\]
		where Precision = TP/(TP+FP)
		
		\vspace{0.1cm}
		\textbf{INTUITIVE APPEAL:}
		\begin{itemize}
			\item Avoid stark differences in false positives across groups
			\item Ensure proportion of accepted applicants who succeed is same
			\item Prevent academic potential being overestimated in one group
		\end{itemize}
		
		\textcolor{myred}{\textbf{EXAM TRAP:}} Predictive rate parity = \textcolor{myblue}{equal precision}!
	\end{frame}
	
	\begin{frame}{T14: Predictive Rate Parity - Drawbacks}
		\fontsize{6.5}{7.5}\selectfont
		\textbf{TOPIC: Drawbacks of Predictive Rate Parity}
		
		\vspace{0.1cm}
		\textbf{MAJOR DRAWBACK:}
		\begin{itemize}
			\item Does not account for differences in base rates across groups
			\item If base rates differ significantly, achieving parity is challenging
			\item Can have adverse effects on algorithmic performance
		\end{itemize}
		
		\vspace{0.1cm}
		\textbf{WHY CHALLENGING:}
		\begin{itemize}
			\item Base rates = general prevalence within population
			\item Disease prevalence may differ across groups
			\item Default rates may differ across demographics
			\item Achieving parity may require sacrificing accuracy
		\end{itemize}
		
		\vspace{0.1cm}
		\textbf{EXAMPLES:}
		\begin{itemize}
			\item Medical diagnostics: disease rates differ
			\item Credit scoring: default rates differ
			\item College admissions: qualification rates differ
		\end{itemize}
		
		\textcolor{myred}{\textbf{EXAM TRAP:}} Predictive rate parity struggles with \textcolor{myblue}{different base rates}!
	\end{frame}
	
	\begin{frame}{T15: Equal Opportunity}
		\fontsize{6.5}{7.5}\selectfont
		\textbf{TOPIC: Equal Opportunity - Definition and Formula}
		
		\vspace{0.1cm}
		\textbf{DEFINITION:}
		\begin{itemize}
			\item Sensitivity (True Positive Rate) is equal across groups
			\item Individuals of each group have equal chance of being correctly identified as positive
			\item Important when positive cases must be recognized
		\end{itemize}
		
		\vspace{0.1cm}
		\textbf{FORMULA:}
		\[
		P(\hat{Y}=1|Y=1,A=0) = P(\hat{Y}=1|Y=1,A=1)
		\]
		
		\vspace{0.1cm}
		\textbf{IN TERMS OF SENSITIVITY:}
		\[
		\text{Sensitivity}_0 = \text{Sensitivity}_1
		\]
		where Sensitivity = TP/(TP+FN)
		
		\vspace{0.1cm}
		\textbf{EXAMPLE (Medical Diagnostics):}
		\begin{itemize}
			\item Patients with disease equally likely to be diagnosed correctly
			\item Regardless of demographic group
			\item Don't miss disease in any group
		\end{itemize}
		
		\textcolor{myred}{\textbf{EXAM TRAP:}} Equal opportunity = \textcolor{myblue}{equal true positive rate (sensitivity)}!
	\end{frame}
	
	\begin{frame}{T16: Equal Opportunity - Drawbacks}
		\fontsize{6.5}{7.5}\selectfont
		\textbf{TOPIC: Drawbacks of Equal Opportunity}
		
		\vspace{0.1cm}
		\textbf{MAJOR DRAWBACK:}
		\begin{itemize}
			\item Requires access to true value of target
			\item Can be difficult to achieve with stark base rate differences
			\item Difficult to achieve fairness without compromising other performance aspects
		\end{itemize}
		
		\vspace{0.1cm}
		\textbf{RELATED CONCEPT - EQUALIZED ODDS:}
		\begin{itemize}
			\item More restrictive version of equal opportunity
			\item Requires equality of sensitivity AND specificity across groups
			\item $P(\hat{Y}=1|Y=1,A=0) = P(\hat{Y}=1|Y=1,A=1)$
			\item $P(\hat{Y}=1|Y=0,A=0) = P(\hat{Y}=1|Y=0,A=1)$
		\end{itemize}
		
		\vspace{0.1cm}
		\textbf{WHEN USEFUL:}
		\begin{itemize}
			\item Ensure same proportion of qualified candidates admitted
			\item Ensure don't miss disease in any group
		\end{itemize}
		
		\textcolor{myred}{\textbf{EXAM TRAP:}} Equal opportunity requires \textcolor{myblue}{access to true values}!
	\end{frame}
	
	\begin{frame}{T17: Fairness Impossibility Theorem}
		\fontsize{6.5}{7.5}\selectfont
		\textbf{TOPIC: Fairness Impossibility Theorem and Trade-offs}
		
		\vspace{0.1cm}
		\textbf{DEFINITION:}
		\begin{itemize}
			\item Mathematical theorem demonstrates impossibility
			\item When base rates differ between populations (almost always)
			\item Cannot satisfy demographic parity, predictive rate parity, and equal opportunity simultaneously
		\end{itemize}
		
		\vspace{0.1cm}
		\textbf{IMPLICATIONS:}
		\begin{itemize}
			\item Cannot achieve fairness on all dimensions
			\item Must deliberate on kind of group fairness to achieve
			\item Trade-offs are unavoidable
			\item Must choose based on context
		\end{itemize}
		
		\vspace{0.1cm}
		\textbf{PERFORMANCE TRADE-OFFS:}
		\begin{itemize}
			\item Balancing different performance measures is non-trivial
			\item Example: Fraud detection - sensitivity vs. false-positive rate
			\item Block all transactions: maximum sensitivity, high false-positive rate
			\item Label all as non-fraudulent: minimum false-positive rate, low sensitivity
		\end{itemize}
		
		\textcolor{myred}{\textbf{EXAM TRAP:}} Cannot satisfy \textcolor{myblue}{all fairness measures} simultaneously!
	\end{frame}
	
	\begin{frame}{T18: Fairness Measures Summary}
		\fontsize{6.5}{7.5}\selectfont
		\textbf{TOPIC: Summary of Group Fairness Measures}
		
		\vspace{0.1cm}
		\textbf{THREE KEY MEASURES:}
		
		\begin{tabular}{|p{2.5cm}|p{5cm}|}
			\hline
			\textbf{Measure} & \textbf{Question Answered} \\
			\hline
			Demographic Parity & "Is the likelihood of a positive prediction the same across groups?" \\
			\hline
			Predictive Rate Parity & "Is the likelihood of a true positive prediction the same across groups?" \\
			\hline
			Equal Opportunity & "Is the likelihood of identifying actual positives as positive the same across groups?" \\
			\hline
		\end{tabular}
		
		\vspace{0.1cm}
		\textbf{KEY CHARACTERISTICS:}
		\begin{itemize}
			\item \textcolor{myblue}{Demographic Parity:} Independent of correctness
			\item \textcolor{myblue}{Predictive Rate Parity:} Focuses on precision
			\item \textcolor{myblue}{Equal Opportunity:} Focuses on sensitivity
		\end{itemize}
		
		\textcolor{myred}{\textbf{EXAM TRAP:}} Know the \textcolor{myblue}{question each measure answers}!
	\end{frame}
	
	\begin{frame}{T19: Performance Metrics - Precision}
		\fontsize{6.5}{7.5}\selectfont
		\textbf{TOPIC: Precision in Classification}
		
		\vspace{0.1cm}
		\textbf{DEFINITION:}
		\begin{itemize}
			\item Proportion of positive predictions that are correct
			\item Also known as Positive Predictive Value (PPV)
		\end{itemize}
		
		\vspace{0.1cm}
		\textbf{FORMULA:}
		\[
		\text{Precision} = \frac{TP}{TP+FP}
		\]
		
		\vspace{0.1cm}
		\textbf{INTERPRETATION:}
		\begin{itemize}
			\item "Of those predicted positive, how many were actually positive?"
			\item High precision: Few false positives
			\item Low precision: Many false positives
		\end{itemize}
		
		\vspace{0.1cm}
		\textbf{EXAMPLE (Credit Scoring):}
		\begin{itemize}
			\item Proportion of applicants predicted to default who actually default
			\item High precision: Algorithm effectively identifies likely defaulters
			\item Low precision: Many who would repay erroneously labeled high-risk
		\end{itemize}
		
		\textcolor{myred}{\textbf{EXAM TRAP:}} Precision = \textcolor{myblue}{TP/(TP+FP)}!
	\end{frame}
	
	\begin{frame}{T20: Performance Metrics - Sensitivity}
		\fontsize{6.5}{7.5}\selectfont
		\textbf{TOPIC: Sensitivity (Recall) in Classification}
		
		\vspace{0.1cm}
		\textbf{DEFINITION:}
		\begin{itemize}
			\item Proportion of actual positive instances correctly predicted
			\item Also known as True Positive Rate (TPR) or Recall
		\end{itemize}
		
		\vspace{0.1cm}
		\textbf{FORMULA:}
		\[
		\text{Sensitivity} = \frac{TP}{TP+FN}
		\]
		
		\vspace{0.1cm}
		\textbf{INTERPRETATION:}
		\begin{itemize}
			\item "Of those actually positive, how many did we catch?"
			\item Useful when important not to miss a positive
			\item Maximizing sensitivity requires minimizing false negatives
		\end{itemize}
		
		\vspace{0.1cm}
		\textbf{RELATED CONCEPT:}
		\begin{itemize}
			\item \textcolor{myblue}{Specificity:} Proportion of actual negatives identified
			\item $Specificity = \frac{TN}{TN+FP}$
		\end{itemize}
		
		\textcolor{myred}{\textbf{EXAM TRAP:}} Sensitivity = \textcolor{myblue}{TP/(TP+FN)}!
	\end{frame}
	
	\begin{frame}{T21: Performance Metrics - Accuracy}
		\fontsize{6.5}{7.5}\selectfont
		\textbf{TOPIC: Accuracy in Classification}
		
		\vspace{0.1cm}
		\textbf{DEFINITION:}
		\begin{itemize}
			\item Fraction of predictions that are correct
			\item Intuitive but often misleading
		\end{itemize}
		
		\vspace{0.1cm}
		\textbf{FORMULA:}
		\[
		\text{Accuracy} = \frac{TP+TN}{TP+TN+FP+FN}
		\]
		
		\vspace{0.1cm}
		\textbf{PROBLEM WITH ACCURACY:}
		\begin{itemize}
			\item Misleading with unbalanced data
			\item Example: 100 applicants, 10 suitable, 90 unsuitable
			\item Algorithm classifies all as unsuitable: Accuracy = 90\%
			\item But provides no useful information
		\end{itemize}
		
		\vspace{0.1cm}
		\textbf{WHEN TO USE:}
		\begin{itemize}
			\item Only when classes are balanced
			\item Or when costs of errors are equal
		\end{itemize}
		
		\textcolor{myred}{\textbf{EXAM TRAP:}} Accuracy is \textcolor{myblue}{misleading with unbalanced data}!
	\end{frame}
	
	\begin{frame}{T22: Performance Metrics - False Positive Rate}
		\fontsize{6.5}{7.5}\selectfont
		\textbf{TOPIC: False Positive Rate (FPR)}
		
		\vspace{0.1cm}
		\textbf{DEFINITION:}
		\begin{itemize}
			\item Proportion of false positives out of all actual negative cases
		\end{itemize}
		
		\vspace{0.1cm}
		\textbf{FORMULA:}
		\[
		FPR = \frac{FP}{FP+TN}
		\]
		
		\vspace{0.1cm}
		\textbf{RELATIONSHIP TO SPECIFICITY:}
		\[
		FPR = 1 - \text{Specificity}
		\]
		
		\vspace{0.1cm}
		\textbf{INTERPRETATION:}
		\begin{itemize}
			\item "Of all actual negatives, how many were incorrectly flagged positive?"
			\item Important for fairness analysis
			\item Used in ROC curves
		\end{itemize}
		
		\textcolor{myred}{\textbf{EXAM TRAP:}} FPR = \textcolor{myblue}{FP/(FP+TN)}!
	\end{frame}
	
	\begin{frame}{T23: Performance Metrics - Confusion Matrix}
		\fontsize{6.5}{7.5}\selectfont
		\textbf{TOPIC: Confusion Matrix}
		
		\vspace{0.1cm}
		\textbf{DEFINITION:}
		\begin{itemize}
			\item Cross-tabulation of predicted vs. actual outcomes
			\item Foundation for all classification metrics
		\end{itemize}
		
		\vspace{0.1cm}
		\textbf{STRUCTURE:}
		
		\begin{tabular}{|c|c|c|}
			\hline
			& \textbf{Predicted Neg} & \textbf{Predicted Pos} \\
			\hline
			\textbf{Actual Neg} & TN & FP \\
			\hline
			\textbf{Actual Pos} & FN & TP \\
			\hline
		\end{tabular}
		
		\vspace{0.1cm}
		\textbf{COMPONENTS:}
		\begin{itemize}
			\item \textcolor{myblue}{True Negative (TN):} Correctly predicted negative
			\item \textcolor{myblue}{False Positive (FP):} Incorrectly predicted positive
			\item \textcolor{myblue}{False Negative (FN):} Incorrectly predicted negative
			\item \textcolor{myblue}{True Positive (TP):} Correctly predicted positive
		\end{itemize}
		
		\textcolor{myred}{\textbf{EXAM TRAP:}} Confusion matrix = \textcolor{myblue}{TP, TN, FP, FN}!
	\end{frame}
	
	% ============================================================
	% SECTION 3: SOURCES OF UNFAIRNESS
	% ============================================================
	
	\begin{frame}{T24: Problem Specification}
		\fontsize{6.5}{7.5}\selectfont
		\textbf{TOPIC: Problem Specification and Feature Selection}
		
		\vspace{0.1cm}
		\textbf{DEFINITION:}
		\begin{itemize}
			\item During problem specification, purpose and design of algorithm are explicitly stated
			\item How a problem is specified affects algorithmic performance and fairness
		\end{itemize}
		
		\vspace{0.1cm}
		\textbf{KEY ISSUES:}
		\begin{itemize}
			\item Choice of criteria affects who applies and who gets hired
			\item Domain expertise is necessary for good problem specification
			\item Translation to measurable features requires proxies
			\item Some things harder to quantify than others
		\end{itemize}
		
		\vspace{0.1cm}
		\textbf{EXAMPLE (Hiring Algorithm):}
		\begin{itemize}
			\item Company wants to improve recruitment using HireMe
			\item Scans CVs, outputs shortlist of best applicants
			\item Need to specify what is a "good applicant"
			\item Productivity, work ethic, team player, retention
		\end{itemize}
		
		\textcolor{myred}{\textbf{EXAM TRAP:}} Problem specification = \textcolor{myblue}{critical for fairness}!
	\end{frame}
	
	\begin{frame}{T25: Proxy Variables}
		\fontsize{6.5}{7.5}\selectfont
		\textbf{TOPIC: Proxy Variables in Algorithmic Design}
		
		\vspace{0.1cm}
		\textbf{DEFINITION:}
		\begin{itemize}
			\item Variables used as stand-ins for concepts hard to measure directly
			\item Judged to be strongly correlated with the criterion
		\end{itemize}
		
		\vspace{0.1cm}
		\textbf{WHY PROXIES ARE USED:}
		\begin{itemize}
			\item Cannot easily measure productivity directly
			\item Can measure sales made within one year
			\item Simpler measurable construct
		\end{itemize}
		
		\vspace{0.1cm}
		\textbf{RISKS:}
		\begin{itemize}
			\item Reduces complex concepts
			\item May introduce bias
			\item Example: Total sales misses that women more likely to work part-time
			\item Lower total sales rate for women
		\end{itemize}
		
		\vspace{0.1cm}
		\textbf{EXAMPLE:}
		\begin{itemize}
			\item Productivity measured by sales
			\item Teamwork harder to quantify
			\item Skipping features may hurt candidates
			\item If teamwork area where women outperform men
		\end{itemize}
		
		\textcolor{myred}{\textbf{EXAM TRAP:}} Proxies are \textcolor{myblue}{imperfect stand-ins}!
	\end{frame}
	
	\begin{frame}{T26: Amazon Recruiting Tool}
		\fontsize{6.5}{7.5}\selectfont
		\textbf{TOPIC: Amazon Recruiting Tool Example}
		
		\vspace{0.1cm}
		\textbf{SCENARIO:}
		\begin{itemize}
			\item Amazon attempted to create unbiased recruiting tool
			\item Provided algorithm with previous applicants' CVs
			\item Target variable: whether hired
			\item "Good" candidate = resembles previous hires
		\end{itemize}
		
		\vspace{0.1cm}
		\textbf{OUTCOME:}
		\begin{itemize}
			\item Algorithm systematically downgraded women's CVs
			\item Even identical CVs rated lower for women
			\item Amazon abandoned the project
		\end{itemize}
		
		\vspace{0.1cm}
		\textbf{LESSON:}
		\begin{itemize}
			\item Learning from historical data perpetuates biases
			\item Past discrimination reflected in training data
			\item Cannot simply rely on past patterns
		\end{itemize}
		
		\textcolor{myred}{\textbf{EXAM TRAP:}} Amazon's tool \textcolor{myblue}{systematically discriminated against women}!
	\end{frame}
	
	\begin{frame}{T27: Fairness Through Unawareness}
		\fontsize{6.5}{7.5}\selectfont
		\textbf{TOPIC: Fairness Through Unawareness - Does It Work?}
		
		\vspace{0.1cm}
		\textbf{DEFINITION:}
		\begin{itemize}
			\item Intuitive approach: Remove protected characteristics from inputs
			\item Assumes this prevents discrimination
			\item Embedded in most legal doctrines
		\end{itemize}
		
		\vspace{0.1cm}
		\textbf{PROBLEM:}
		\begin{itemize}
			\item Correlated variables carry protected information
			\item Removing protected variable doesn't eliminate its influence
			\item Correlated variables weighted more strongly
		\end{itemize}
		
		\vspace{0.1cm}
		\textbf{EXAMPLE (FTA Prediction):}
		\begin{itemize}
			\item JusticeGuard predicts failure to appear in court
			\item Protected group and prior convictions correlated
			\item Remove protected variable
			\item Prior conviction coefficient increases from 54\% to 72\%
		\end{itemize}
		
		\textcolor{myred}{\textbf{EXAM TRAP:}} Fairness through unawareness \textcolor{myblue}{does NOT eliminate bias}!
	\end{frame}
	
	\begin{frame}{T28: FTA Prediction Example}
		\fontsize{6.5}{7.5}\selectfont
		\textbf{TOPIC: FTA Prediction - Mathematical Illustration}
		
		\vspace{0.1cm}
		\textbf{INITIAL MODEL (with protected characteristic):}
		\[
		\hat{Y} = 0.034 + 0.543 X_{prior} + 0.332 X_{prot}
		\]
		\begin{itemize}
			\item Prior offense: 54\% increase in risk
			\item Protected group: 33\% increase in risk
		\end{itemize}
		
		\vspace{0.1cm}
		\textbf{REMOVE PROTECTED CHARACTERISTIC:}
		\[
		\hat{Y} = \gamma_0 + \gamma_1 X_{prior}
		\]
		
		\vspace{0.1cm}
		\textbf{RESULT:}
		\[
		\hat{Y} = 0.10 + 0.72 X_{prior}
		\]
		\begin{itemize}
			\item Prior now increases risk by 72\% (not 54\%)
			\item Coefficient contaminated by correlation
			\item Algorithm still indirectly discriminates
		\end{itemize}
		
		\textcolor{myred}{\textbf{EXAM TRAP:}} Removing protected variable \textcolor{myblue}{can make bias worse}!
	\end{frame}
	
	\begin{frame}{T29: Historical Bias}
		\fontsize{6.5}{7.5}\selectfont
		\textbf{TOPIC: Historical Bias in Training Data}
		
		\vspace{0.1cm}
		\textbf{DEFINITION:}
		\begin{itemize}
			\item Supervised learning relies on historical data
			\item Historical data reflects past biases and discrimination
			\item Algorithm learns and may exacerbate these patterns
		\end{itemize}
		
		\vspace{0.1cm}
		\textbf{EXAMPLE:}
		\begin{itemize}
			\item Hiring algorithm trained on past hiring data
			\item Women historically less likely selected for certain roles
			\item Algorithm repeats same patterns
			\item May exacerbate them
		\end{itemize}
		
		\vspace{0.1cm}
		\textbf{CONSEQUENCE:}
		\begin{itemize}
			\item Systematic bias perpetuated
			\item Discrimination becomes automated
			\item Hard to detect and correct
		\end{itemize}
		
		\textcolor{myred}{\textbf{EXAM TRAP:}} Historical bias = \textcolor{myblue}{past discrimination reflected in data}!
	\end{frame}
	
	\begin{frame}{T30: Feedback Loops}
		\fontsize{6.5}{7.5}\selectfont
		\textbf{TOPIC: Feedback Loops in AI Systems}
		
		\vspace{0.1cm}
		\textbf{DEFINITION:}
		\begin{itemize}
			\item Algorithmic predictions influence future data
			\item Biased predictions create more biased data
			\item Amplifies existing discrimination
		\end{itemize}
		
		\vspace{0.1cm}
		\textbf{HOW IT WORKS:}
		\begin{enumerate}
			\item Algorithm makes biased predictions
			\item Decisions based on predictions create new data
			\item New data reflects biased decisions
			\item Algorithm learns from biased data
			\item Bias amplified
		\end{enumerate}
		
		\vspace{0.1cm}
		\textbf{EXAMPLE:}
		\begin{itemize}
			\item Hiring algorithm learns from biased outcomes
			\item Fewer women hired
			\item Algorithm sees fewer women in successful hires
			\item Future predictions even more biased
		\end{itemize}
		
		\textcolor{myred}{\textbf{EXAM TRAP:}} Feedback loops \textcolor{myblue}{amplify and perpetuate bias}!
	\end{frame}
	
	\begin{frame}{T31: Representation Bias}
		\fontsize{6.5}{7.5}\selectfont
		\textbf{TOPIC: Representation Bias in Data Collection}
		
		\vspace{0.1cm}
		\textbf{DEFINITION:}
		\begin{itemize}
			\item Collected sample not representative of target population
			\item Results in problematic model specifications
		\end{itemize}
		
		\vspace{0.1cm}
		\textbf{EXAMPLES:}
		\begin{itemize}
			\item Facial recognition trained mostly on white faces
			\item Medical testing with predominantly white/European samples
			\item Algorithms perform better on dominant groups
		\end{itemize}
		
		\vspace{0.1cm}
		\textbf{INTERSECTIONAL BIAS:}
		\begin{itemize}
			\item Performance particularly low for intersectional groups
			\item Example: Women of color in facial recognition
			\item Multiple dimensions of underrepresentation
		\end{itemize}
		
		\textcolor{myred}{\textbf{EXAM TRAP:}} Representation bias = \textcolor{myblue}{non-representative sample}!
	\end{frame}
	
	\begin{frame}{T32: Measurement Bias}
		\fontsize{6.5}{7.5}\selectfont
		\textbf{TOPIC: Measurement Bias in Data Collection}
		
		\vspace{0.1cm}
		\textbf{DEFINITION:}
		\begin{itemize}
			\item Measurement methods not consistent across sample
			\item Can favor certain groups
		\end{itemize}
		
		\vspace{0.1cm}
		\textbf{EXAMPLE:}
		\begin{itemize}
			\item Doctor spends more time with wealthy patients
			\item More thorough evaluations
			\item Data on early-stage disease discovery biased
			\item Favor of wealthy group
		\end{itemize}
		
		\vspace{0.1cm}
		\textbf{CONSEQUENCE:}
		\begin{itemize}
			\item Biased training data
			\item Algorithm learns incorrect patterns
			\item Systematic errors for certain groups
		\end{itemize}
		
		\textcolor{myred}{\textbf{EXAM TRAP:}} Measurement bias = \textcolor{myblue}{inconsistent measurement across groups}!
	\end{frame}
	
	\begin{frame}{T33: Data Composition Bias}
		\fontsize{6.5}{7.5}\selectfont
		\textbf{TOPIC: Data Composition Bias}
		
		\vspace{0.1cm}
		\textbf{DEFINITION:}
		\begin{itemize}
			\item Minority subgroups underrepresented in data
			\item Algorithms underperform for these groups
			\item Statistical nature of learning delivers statistically ideal output
		\end{itemize}
		
		\vspace{0.1cm}
		\textbf{EXAMPLE:}
		\begin{itemize}
			\item Medical training set: 3\% pregnant women
			\item Pregnant women require different treatment
			\item Algorithm significantly underperforms for this group
			\item Fewer data points to learn from
		\end{itemize}
		
		\vspace{0.1cm}
		\textbf{SOLUTION:}
		\begin{itemize}
			\item Balance datasets
			\item Relevant minority subgroups occupy similar proportion
			\item Ensures algorithm recognizes minority groups as statistically relevant
		\end{itemize}
		
		\textcolor{myred}{\textbf{EXAM TRAP:}} Data composition bias = \textcolor{myblue}{underrepresented minority subgroups}!
	\end{frame}
	
	\begin{frame}{T34: Balancing Datasets}
		\fontsize{6.5}{7.5}\selectfont
		\textbf{TOPIC: Balancing Datasets}
		
		\vspace{0.1cm}
		\textbf{DEFINITION:}
		\begin{itemize}
			\item Ensuring relevant minority subgroups occupy similar proportion as majority groups
			\item Different from proportional representation
		\end{itemize}
		
		\vspace{0.1cm}
		\textbf{EXAMPLE:}
		\begin{itemize}
			\item \textcolor{myblue}{Proportional:} 10\% triangles, 90\% circles → dataset same
			\item \textcolor{myblue}{Balanced:} Increase proportion of triangles in dataset
			\item Male/female 50\% even if historically 20\% female
		\end{itemize}
		
		\vspace{0.1cm}
		\textbf{METHODS:}
		\begin{itemize}
			\item \textcolor{myblue}{Over-sampling:} Add more data from minority group
			\item \textcolor{myblue}{Under-sampling:} Remove data from majority group
			\item \textcolor{myblue}{Combination:} Both methods
			\item \textcolor{myblue}{Synthetic data:} Generate new data points
		\end{itemize}
		
		\textcolor{myred}{\textbf{EXAM TRAP:}} Balancing = \textcolor{myblue}{over-representing minorities} for algorithm performance!
	\end{frame}
	
	\begin{frame}{T35: Bootstrapping Methods}
		\fontsize{6.5}{7.5}\selectfont
		\textbf{TOPIC: Bootstrapping for Data Balancing}
		
		\vspace{0.1cm}
		\textbf{DEFINITION:}
		\begin{itemize}
			\item Resampling datapoints from existing dataset
			\item Adding back into dataset
		\end{itemize}
		
		\vspace{0.1cm}
		\textbf{THREE METHODS:}
		\begin{enumerate}
			\item \textcolor{myblue}{Over-sampling:} Datapoints from minority group
			\item \textcolor{myblue}{Under-sampling:} Remove datapoints from majority group
			\item \textcolor{myblue}{Combination:} Both methods
		\end{enumerate}
		
		\vspace{0.1cm}
		\textbf{CHALLENGES:}
		\begin{itemize}
			\item Creating synthetic data without compromising generalization
			\item Taking out datapoints works only if datasets are large enough
			\item Rare cases especially challenging (fraud <1\%)
		\end{itemize}
		
		\vspace{0.1cm}
		\textbf{ADDITIONAL METHODS:}
		\begin{itemize}
			\item Ensemble methods: Use multiple models
			\item Cost function: Higher penalty for minority class misclassification
		\end{itemize}
		
		\textcolor{myred}{\textbf{EXAM TRAP:}} Bootstrapping includes \textcolor{myblue}{over-sampling AND under-sampling}!
	\end{frame}
	
	\begin{frame}{T36: Intersectional Bias}
		\fontsize{6.5}{7.5}\selectfont
		\textbf{TOPIC: Intersectional Bias}
		
		\vspace{0.1cm}
		\textbf{DEFINITION:}
		\begin{itemize}
			\item Bias affecting groups at intersection of multiple underrepresented subgroups
			\item Performance particularly low for intersectional groups
		\end{itemize}
		
		\vspace{0.1cm}
		\textbf{EXAMPLE:}
		\begin{itemize}
			\item Women of color in facial recognition
			\item Multiple dimensions of underrepresentation
			\item Harder to detect and correct
		\end{itemize}
		
		\vspace{0.1cm}
		\textbf{WHY IT MATTERS:}
		\begin{itemize}
			\item Standard fairness metrics may miss intersectional bias
			\item Need to analyze multiple dimensions
			\item Particularly harmful for affected individuals
		\end{itemize}
		
		\textcolor{myred}{\textbf{EXAM TRAP:}} Intersectional bias = \textcolor{myblue}{multiple underrepresented dimensions combined}!
	\end{frame}
	
	\begin{frame}{T37: Word Embeddings Bias}
		\fontsize{6.5}{7.5}\selectfont
		\textbf{TOPIC: Word Embeddings Bias}
		
		\vspace{0.1cm}
		\textbf{DEFINITION:}
		\begin{itemize}
			\item Preprocessing technique for NLP
			\item Transforms words into high-dimensional vectors
			\item Captures meaning through relationships to other words
		\end{itemize}
		
		\vspace{0.1cm}
		\textbf{EXAMPLES:}
		\begin{itemize}
			\item man is to woman as king is to queen
			\item man is to computer programmer as woman is to homemaker
			\item Contains stereotypes
		\end{itemize}
		
		\vspace{0.1cm}
		\textbf{CONSEQUENCE:}
		\begin{itemize}
			\item NLP algorithms incorporate stereotypes
			\item May perpetuate harmful biases
			\item Difficult to detect and correct
		\end{itemize}
		
		\textcolor{myred}{\textbf{EXAM TRAP:}} Word embeddings can contain \textcolor{myblue}{harmful stereotypes}!
	\end{frame}
	
	\begin{frame}{T38: Objective Function Bias}
		\fontsize{6.5}{7.5}\selectfont
		\textbf{TOPIC: Objective Function Bias}
		
		\vspace{0.1cm}
		\textbf{DEFINITION:}
		\begin{itemize}
			\item Choice of objective function introduces value judgments
			\item What algorithm optimizes for makes ethically significant difference
		\end{itemize}
		
		\vspace{0.1cm}
		\textbf{KEY QUESTIONS:}
		\begin{itemize}
			\item Minimize prediction errors (focus on performance)?
			\item Or focus on how errors are distributed (focus on error distribution)?
			\item Choice depends on context
		\end{itemize}
		
		\vspace{0.1cm}
		\textbf{EXAMPLE (Cancer Detection):}
		\begin{itemize}
			\item \textcolor{myblue}{Skin cancer:} Err on side of caution
			\item \textcolor{myblue}{Brain cancer:} Treatment highly invasive
			\item Different trade-offs for false positives/negatives
			\item Context matters for objective function
		\end{itemize}
		
		\textcolor{myred}{\textbf{EXAM TRAP:}} Objective function choice = \textcolor{myblue}{ethical significance}!
	\end{frame}
	
	\begin{frame}{T39: Benchmark Dataset Bias}
		\fontsize{6.5}{7.5}\selectfont
		\textbf{TOPIC: Benchmark Dataset Bias}
		
		\vspace{0.1cm}
		\textbf{DEFINITION:}
		\begin{itemize}
			\item Algorithms tested against publicly available benchmark datasets
			\item May not be representative of intended use population
			\item Not necessarily good indicator of model quality
		\end{itemize}
		
		\vspace{0.1cm}
		\textbf{PROBLEM:}
		\begin{itemize}
			\item Good performance on benchmark doesn't guarantee real-world performance
			\item Can mask biases that appear in deployment
			\item Need domain-specific testing
		\end{itemize}
		
		\textcolor{myred}{\textbf{EXAM TRAP:}} Benchmark datasets may \textcolor{myblue}{not reflect real-world populations}!
	\end{frame}
	
	\begin{frame}{T40: Deployment Context Bias}
		\fontsize{6.5}{7.5}\selectfont
		\textbf{TOPIC: Deployment Context Bias}
		
		\vspace{0.1cm}
		\textbf{DEFINITION:}
		\begin{itemize}
			\item Algorithms used in different contexts than designed for
			\item Loss of performance or other challenges
			\item Applies to different contexts and outdated data
		\end{itemize}
		
		\vspace{0.1cm}
		\textbf{EXAMPLES:}
		\begin{itemize}
			\item Recidivism prediction for sentence length
			\item Values of users differ from developers
			\item Third-party algorithms may not align with company values
		\end{itemize}
		
		\vspace{0.1cm}
		\textbf{SOLUTIONS:}
		\begin{itemize}
			\item Human oversight of data collection
			\item Oversight of development
			\item Oversight of deployment
			\item Regular auditing
		\end{itemize}
		
		\textcolor{myred}{\textbf{EXAM TRAP:}} Deployment context = \textcolor{myblue}{critical for performance and fairness}!
	\end{frame}
	
	\begin{frame}{T41: Human Oversight}
		\fontsize{6.5}{7.5}\selectfont
		\textbf{TOPIC: Importance of Human Oversight}
		
		\vspace{0.1cm}
		\textbf{DEFINITION:}
		\begin{itemize}
			\item Ensuring safe and fair algorithms requires human oversight
			\item Throughout AI lifecycle
		\end{itemize}
		
		\vspace{0.1cm}
		\textbf{AREAS REQUIRING OVERSIGHT:}
		\begin{itemize}
			\item Data collection process
			\item Development of algorithm
			\item Deployment
			\item Regular auditing
		\end{itemize}
		
		\vspace{0.1cm}
		\textbf{WHY IT MATTERS:}
		\begin{itemize}
			\item Algorithms deployed in consequential domains
			\item Need accountability
			\item Regular auditing becomes important
			\item Failures can have severe consequences
		\end{itemize}
		
		\textcolor{myred}{\textbf{EXAM TRAP:}} Human oversight is \textcolor{myblue}{essential throughout AI lifecycle}!
	\end{frame}
	
	\begin{frame}{T42: Sources of Unfairness Summary}
		\fontsize{6.5}{7.5}\selectfont
		\textbf{TOPIC: Summary of Sources of Unfairness}
		
		\vspace{0.1cm}
		\textbf{MAJOR SOURCES:}
		\begin{enumerate}
			\item \textcolor{myblue}{Problem Specification:}
			\begin{itemize}
				\item Choice of criteria
				\item Translation to measurable features
				\item Proxies
			\end{itemize}
			\item \textcolor{myblue}{Data Collection:}
			\begin{itemize}
				\item Historical bias
				\item Representation bias
				\item Measurement bias
			\end{itemize}
			\item \textcolor{myblue}{Data Composition:}
			\begin{itemize}
				\item Imbalanced datasets
				\item Statistical minorities
			\end{itemize}
			\item \textcolor{myblue}{Model Development:}
			\begin{itemize}
				\item Preprocessing techniques
				\item Choice of objective function
				\item Benchmark testing
			\end{itemize}
			\item \textcolor{myblue}{Deployment:}
			\begin{itemize}
				\item Context mismatch
				\item Value misalignment
			\end{itemize}
		\end{enumerate}
		
		\textcolor{myred}{\textbf{EXAM TRAP:}} Multiple sources - need \textcolor{myblue}{holistic approach}!
	\end{frame}
	
	% ============================================================
	% SECTION 4: EXPLAINABILITY, INTERPRETABILITY, TRANSPARENCY
	% ============================================================
	
	\begin{frame}{T43: Explainability}
		\fontsize{6.5}{7.5}\selectfont
		\textbf{TOPIC: Explainability in AI}
		
		\vspace{0.1cm}
		\textbf{DEFINITION:}
		\begin{itemize}
			\item Capacity to explain in understandable terms how AI system makes decisions
			\item Often after the fact (ex-post explainability)
		\end{itemize}
		
		\vspace{0.1cm}
		\textbf{WHY IT MATTERS:}
		\begin{itemize}
			\item Inability to explain leads to accountability gaps
			\item Erosion of trust among users and stakeholders
			\item Essential in consequential domains (healthcare, finance)
		\end{itemize}
		
		\vspace{0.1cm}
		\textbf{TWO ASPECTS:}
		\begin{itemize}
			\item \textcolor{myblue}{Philosophical:} What counts as a good explanation?
			\item \textcolor{myblue}{Technical:} How to obtain good explanations?
		\end{itemize}
		
		\textcolor{myred}{\textbf{EXAM TRAP:}} Explainability = \textcolor{myblue}{explaining decisions after the fact}!
	\end{frame}
	
	\begin{frame}{T44: Interpretability}
		\fontsize{6.5}{7.5}\selectfont
		\textbf{TOPIC: Interpretability in AI}
		
		\vspace{0.1cm}
		\textbf{DEFINITION:}
		\begin{itemize}
			\item Degree to which human can comprehend and predict model's behavior
			\item Built-in mechanisms to understand inherently
			\item How inputs affect outputs
		\end{itemize}
		
		\vspace{0.1cm}
		\textbf{INHERENTLY INTERPRETABLE MODELS:}
		\begin{itemize}
			\item Decision trees: Can follow logic
			\item Linear regression: Coefficients explain impact
			\item Rule-based systems: Explicit rules
		\end{itemize}
		
		\vspace{0.1cm}
		\textbf{BLACK-BOX MODELS:}
		\begin{itemize}
			\item Neural networks: Complex architecture
			\item Ensemble methods: Hard to trace logic
			\item Less interpretable
		\end{itemize}
		
		\textcolor{myred}{\textbf{EXAM TRAP:}} Interpretability = \textcolor{myblue}{inherently understandable models}!
	\end{frame}
	
	\begin{frame}{T45: Transparency}
		\fontsize{6.5}{7.5}\selectfont
		\textbf{TOPIC: Transparency in AI}
		
		\vspace{0.1cm}
		\textbf{DEFINITION:}
		\begin{itemize}
			\item Openness and accessibility of AI system's design, algorithms, data, and decision-making processes
			\item Builds trust, facilitates collaboration
			\item Enables oversight and governance
		\end{itemize}
		
		\vspace{0.1cm}
		\textbf{WHY IT MATTERS:}
		\begin{itemize}
			\item Essential for accountability
			\item Enables alignment with ethical standards
			\item Aligns with organizational values
			\item Meets regulatory requirements
		\end{itemize}
		
		\vspace{0.1cm}
		\textbf{CHALLENGES:}
		\begin{itemize}
			\item Security concerns (spam filters)
			\item Proprietary interests
			\item Specialized knowledge required
		\end{itemize}
		
		\textcolor{myred}{\textbf{EXAM TRAP:}} Transparency = \textcolor{myblue}{openness and accessibility}!
	\end{frame}
	
	\begin{frame}{T46: Black-Box Problem}
		\fontsize{6.5}{7.5}\selectfont
		\textbf{TOPIC: The Black-Box Problem}
		
		\vspace{0.1cm}
		\textbf{DEFINITION:}
		\begin{itemize}
			\item Model too complex to be interpreted
			\item Humans cannot understand intrinsic architecture
			\item Leads to explainability issues
		\end{itemize}
		
		\vspace{0.1cm}
		\textbf{EXAMPLES:}
		\begin{itemize}
			\item Neural networks: Millions of parameters
			\item Deep learning models: Complex interactions
			\item Ensemble methods: Hard to trace
		\end{itemize}
		
		\vspace{0.1cm}
		\textbf{CONSEQUENCES:}
		\begin{itemize}
			\item Understanding decision-making process difficult
			\item If not impossible
			\item Accountability gaps
			\item Trust issues
		\end{itemize}
		
		\textcolor{myred}{\textbf{EXAM TRAP:}} Black-box problem = \textcolor{myblue}{models too complex to interpret}!
	\end{frame}
	
	\begin{frame}{T47: Feature Importance Scores}
		\fontsize{6.5}{7.5}\selectfont
		\textbf{TOPIC: Feature Importance Scores}
		
		\vspace{0.1cm}
		\textbf{DEFINITION:}
		\begin{itemize}
			\item Prominent technique in Explainable AI (XAI)
			\item Ranks importance of input features for outcomes
		\end{itemize}
		
		\vspace{0.1cm}
		\textbf{EXAMPLE:}
		\begin{itemize}
			\item Credit-scoring AI system
			\item Reveals which factors influence credit score
			\item Income, credit history, etc.
		\end{itemize}
		
		\vspace{0.1cm}
		\textbf{BENEFITS:}
		\begin{itemize}
			\item Scrutinize how model makes decisions
			\item Identify potential biases or errors
			\item Analyze impact of input features
		\end{itemize}
		
		\textcolor{myred}{\textbf{EXAM TRAP:}} Feature importance = \textcolor{myblue}{ranks feature influence}!
	\end{frame}
	
	\begin{frame}{T48: Shapley Values}
		\fontsize{6.5}{7.5}\selectfont
		\textbf{TOPIC: Shapley Values in XAI}
		
		\vspace{0.1cm}
		\textbf{DEFINITION:}
		\begin{itemize}
			\item Calculate contribution of each input feature to model output
			\item Measure marginal effect of including each feature
			\item Across all possible feature combinations
		\end{itemize}
		
		\vspace{0.1cm}
		\textbf{BASED ON GAME THEORY:}
		\begin{itemize}
			\item Fair attribution of contributions
			\item Considers all possible coalitions
			\item Provides consistent explanations
		\end{itemize}
		
		\vspace{0.1cm}
		\textbf{APPLICATIONS:}
		\begin{itemize}
			\item Understanding model decisions
			\item Identifying important features
			\item Regulatory explanations
		\end{itemize}
		
		\textcolor{myred}{\textbf{EXAM TRAP:}} Shapley values = \textcolor{myblue}{feature contribution calculation}!
	\end{frame}
	
	\begin{frame}{T49: LUCID}
		\fontsize{6.5}{7.5}\selectfont
		\textbf{TOPIC: LUCID in XAI}
		
		\vspace{0.1cm}
		\textbf{DEFINITION:}
		\begin{itemize}
			\item LUCID = Locating Unfairness through Canonical Inverse Design
			\item Works backward from algorithmic output
			\item Creates distribution over input space
		\end{itemize}
		
		\vspace{0.1cm}
		\textbf{PROCESS:}
		\begin{itemize}
			\item Start from particular output (e.g., "loan granted")
			\item Create distribution over inputs
			\item Reveals model's internal logic
			\item Helps locate unfairness
		\end{itemize}
		
		\vspace{0.1cm}
		\textbf{BENEFITS:}
		\begin{itemize}
			\item Understanding internal logic
			\item Identifying sources of unfairness
			\item Informing mitigation strategies
		\end{itemize}
		
		\textcolor{myred}{\textbf{EXAM TRAP:}} LUCID works \textcolor{myblue}{backward from output to reveal logic}!
	\end{frame}
	
	\begin{frame}{T50: Surrogate Models}
		\fontsize{6.5}{7.5}\selectfont
		\textbf{TOPIC: Surrogate Models in XAI}
		
		\vspace{0.1cm}
		\textbf{DEFINITION:}
		\begin{itemize}
			\item Simpler, more interpretable models
			\item Approximate behavior of complex AI system
		\end{itemize}
		
		\vspace{0.1cm}
		\textbf{EXAMPLE:}
		\begin{itemize}
			\item LIME (Local Interpretable Model-agnostic Explanations)
			\item Approximates complex model's prediction locally
			\item Uses interpretable model (linear regression, decision tree)
			\item Model-agnostic
		\end{itemize}
		
		\vspace{0.1cm}
		\textbf{BENEFITS:}
		\begin{itemize}
			\item Enhance interpretability
			\item Approximate complex model locally
			\item Can be quite accurate locally
		\end{itemize}
		
		\textcolor{myred}{\textbf{EXAM TRAP:}} Surrogate models are \textcolor{myblue}{simpler approximations}!
	\end{frame}
	
	\begin{frame}{T51: XAI Challenges}
		\fontsize{6.5}{7.5}\selectfont
		\textbf{TOPIC: Challenges with Explainable AI}
		
		\vspace{0.1cm}
		\textbf{KEY CHALLENGES:}
		\begin{enumerate}
			\item \textcolor{myblue}{Oversimplification:}
			\begin{itemize}
				\item Complex decisions oversimplified
				\item Can distort understanding
				\item Lead to misinterpretation
			\end{itemize}
			\item \textcolor{myblue}{Computational Cost:}
			\begin{itemize}
				\item Several XAI techniques computationally intensive
				\item More costly to develop and deploy
			\end{itemize}
			\item \textcolor{myblue}{Ex-Post Nature:}
			\begin{itemize}
				\item Used to understand models after training
				\item Not built into model design
			\end{itemize}
		\end{enumerate}
		
		\vspace{0.1cm}
		\textbf{TRADE-OFF:}
		\begin{itemize}
			\item Complex models often achieve higher accuracy
			\item But are less interpretable
			\item Must weigh both factors in context
		\end{itemize}
		
		\textcolor{myred}{\textbf{EXAM TRAP:}} XAI can \textcolor{myblue}{oversimplify and distort understanding}!
	\end{frame}
	
	\begin{frame}{T52: Opaqueness - Secrecy}
		\fontsize{6.5}{7.5}\selectfont
		\textbf{TOPIC: Secrecy as a Form of Opaqueness}
		
		\vspace{0.1cm}
		\textbf{DEFINITION:}
		\begin{itemize}
			\item Individual not aware of being subject to algorithmic decision making
			\item Data extracted without knowledge
		\end{itemize}
		
		\vspace{0.1cm}
		\textbf{EXAMPLE:}
		\begin{itemize}
			\item Algorithm extracts personal data from social media
			\item Clicks, buys, website visits
			\item Extrapolates profile
			\item Makes recommendations for advertisements
			\item Without user's knowledge
		\end{itemize}
		
		\vspace{0.1cm}
		\textbf{VIOLATION:}
		\begin{itemize}
			\item Transparency and consent
			\item User autonomy
			\item Privacy
		\end{itemize}
		
		\textcolor{myred}{\textbf{EXAM TRAP:}} Secrecy = \textcolor{myblue}{unawareness of algorithmic decision making}!
	\end{frame}
	
	\begin{frame}{T53: Opaqueness - Confidentiality}
		\fontsize{6.5}{7.5}\selectfont
		\textbf{TOPIC: Confidentiality as a Form of Opaqueness}
		
		\vspace{0.1cm}
		\textbf{DEFINITION:}
		\begin{itemize}
			\item Affected parties aware of being subject to algorithmic decision making
			\item But don't have knowledge of how decisions are made
			\item Don't understand reasons for outcomes
		\end{itemize}
		
		\vspace{0.1cm}
		\textbf{MAIN CONCERN:}
		\begin{itemize}
			\item Recent discourse on algorithmic decision making
			\item Not unique to algorithms
			\item Secret decision-making in general
		\end{itemize}
		
		\vspace{0.1cm}
		\textbf{EXAMPLE:}
		\begin{itemize}
			\item Automatically denied loan
			\item No way to find out why
			\item Frustration and distrust
		\end{itemize}
		
		\textcolor{myred}{\textbf{EXAM TRAP:}} Confidentiality = \textcolor{myblue}{opaque decision-making process}!
	\end{frame}
	
	\begin{frame}{T54: Opaqueness - Specialized Knowledge}
		\fontsize{6.5}{7.5}\selectfont
		\textbf{TOPIC: Specialized Knowledge as a Form of Opaqueness}
		
		\vspace{0.1cm}
		\textbf{DEFINITION:}
		\begin{itemize}
			\item Understanding algorithm requires specialized training
			\item Even with access to source code and data
			\item May not be able to make sense of it
		\end{itemize}
		
		\vspace{0.1cm}
		\textbf{CHARACTERISTICS:}
		\begin{itemize}
			\item Not unique to AI systems
			\item Results from lack of preexisting knowledge
			\item Creates unequal access to understanding
		\end{itemize}
		
		\vspace{0.1cm}
		\textbf{IMPLICATIONS:}
		\begin{itemize}
			\item Need for expert intermediaries
			\item Challenges for meaningful consent
			\item Accountability gaps
		\end{itemize}
		
		\textcolor{myred}{\textbf{EXAM TRAP:}} Specialized knowledge = \textcolor{myblue}{requires training to understand}!
	\end{frame}
	
	\begin{frame}{T55: Legitimate Reasons for Opacity}
		\fontsize{6.5}{7.5}\selectfont
		\textbf{TOPIC: Valid Reasons for Algorithmic Opacity}
		
		\vspace{0.1cm}
		\textbf{SECURITY:}
		\begin{itemize}
			\item Making code public enables circumvention
			\item Example: Spam filters would be easier to bypass
			\item Risk of fraud, scams, cyber-attacks
		\end{itemize}
		
		\vspace{0.1cm}
		\textbf{PROPRIETARY INTERESTS:}
		\begin{itemize}
			\item Trade secrets
			\item Maintaining competitive advantage
			\item Algorithm as intellectual property
		\end{itemize}
		
		\vspace{0.1cm}
		\textbf{BALANCING ACT:}
		\begin{itemize}
			\item Companies should obtain access for risk assessments
			\item Even with third-party algorithms
			\item Provide accountability
		\end{itemize}
		
		\textcolor{myred}{\textbf{EXAM TRAP:}} Legitimate reasons include \textcolor{myblue}{security and proprietary interests}!
	\end{frame}
	
	% ============================================================
	% SECTION 5: AUTONOMY AND MANIPULATION
	% ============================================================
	
	\begin{frame}{T56: Human Autonomy}
		\fontsize{6.5}{7.5}\selectfont
		\textbf{TOPIC: Human Autonomy in the Age of AI}
		
		\vspace{0.1cm}
		\textbf{DEFINITION:}
		\begin{itemize}
			\item Ability to hold beliefs, make, and execute decisions of our own
			\item Two dimensions
		\end{itemize}
		
		\vspace{0.1cm}
		\textbf{TWO DIMENSIONS:}
		\begin{itemize}
			\item \textcolor{myblue}{Internal:}
			\begin{itemize}
				\item Values, beliefs, decisions that are our own
				\item Not product of distorting external factors
				\item Not manipulation
			\end{itemize}
			\item \textcolor{myblue}{External:}
			\begin{itemize}
				\item Ability to execute decisions
				\item Freedom and opportunity to do so
			\end{itemize}
		\end{itemize}
		
		\vspace{0.1cm}
		\textbf{MISCONCEPTION:}
		\begin{itemize}
			\item Delegating tasks to AI doesn't automatically lose autonomy
			\item Routine outsourcing doesn't reduce autonomy
			\item As long as adequate control and human oversight
		\end{itemize}
		
		\textcolor{myred}{\textbf{EXAM TRAP:}} Autonomy = \textcolor{myblue}{own beliefs and decisions + ability to execute}!
	\end{frame}
	
	\begin{frame}{T57: Cambridge Analytica Scandal}
		\fontsize{6.5}{7.5}\selectfont
		\textbf{TOPIC: Cambridge Analytica Scandal}
		
		\vspace{0.1cm}
		\textbf{SCENARIO:}
		\begin{itemize}
			\item Collected Facebook user data
			\item Users completed personality tests
			\item Inferred personality profiles from online behavior
			\item Sent psychologically tailored political advertising
		\end{itemize}
		
		\vspace{0.1cm}
		\textbf{DEMONSTRATED:}
		\begin{itemize}
			\item Big data can be used to manipulate voter behavior
			\item Potential for political manipulation
			\item Need for robust data privacy and security
		\end{itemize}
		
		\vspace{0.1cm}
		\textbf{REPUTATIONAL IMPACT:}
		\begin{itemize}
			\item Facebook faced reputational damage
			\item Brand value loss
			\item Public trust erosion
		\end{itemize}
		
		\textcolor{myred}{\textbf{EXAM TRAP:}} Cambridge Analytica = \textcolor{myblue}{political manipulation through data}!
	\end{frame}
	
	\begin{frame}{T58: Facebook Emotional Contagion}
		\fontsize{6.5}{7.5}\selectfont
		\textbf{TOPIC: Facebook Emotional Contagion Experiment}
		
		\vspace{0.1cm}
		\textbf{SCENARIO:}
		\begin{itemize}
			\item Researchers changed users' timelines
			\item Displayed predominantly positive or negative posts
			\item Users' own posts became more positive/negative
			\item No consent from users
		\end{itemize}
		
		\vspace{0.1cm}
		\textbf{DEMONSTRATED:}
		\begin{itemize}
			\item Power of recommendation algorithms
			\item Ability to shape experience
			\item Potential to manipulate
		\end{itemize}
		
		\vspace{0.1cm}
		\textbf{PUBLIC OUTCRY:}
		\begin{itemize}
			\item Users did not consent
			\item Violation of trust
			\item Ethical concerns
		\end{itemize}
		
		\textcolor{myred}{\textbf{EXAM TRAP:}} Experiment showed \textcolor{myblue}{AI can manipulate emotional state}!
	\end{frame}
	
	\begin{frame}{T59: AI Manipulation Prevention}
		\fontsize{6.5}{7.5}\selectfont
		\textbf{TOPIC: Preventing AI Manipulation}
		
		\vspace{0.1cm}
		\textbf{APPROACHES:}
		\begin{itemize}
			\item \textcolor{myblue}{Human Oversight:}
			\begin{itemize}
				\item Throughout development process
				\item Meaningful human control
			\end{itemize}
			\item \textcolor{myblue}{Regulatory Oversight:}
			\begin{itemize}
				\item EU AI Act
				\item Transparency obligations
				\item Prohibiting subliminal techniques
			\end{itemize}
		\end{itemize}
		
		\vspace{0.1cm}
		\textbf{CHALLENGES:}
		\begin{itemize}
			\item Ambiguous notion of manipulation
			\item Difficult to operationalize
			\item Hard to make measurable
		\end{itemize}
		
		\vspace{0.1cm}
		\textbf{BALANCE:}
		\begin{itemize}
			\item Benefits of personalization and targeting
			\item vs. risks of manipulation
			\item Need adequate oversight
		\end{itemize}
		
		\textcolor{myred}{\textbf{EXAM TRAP:}} Prevention requires \textcolor{myblue}{human and regulatory oversight}!
	\end{frame}
	
	% ============================================================
	% SECTION 6: SAFETY AND WELL-BEING
	% ============================================================
	
	\begin{frame}{T60: Uber Crash Example}
		\fontsize{6.5}{7.5}\selectfont
		\textbf{TOPIC: Uber Driverless Vehicle Crash}
		
		\vspace{0.1cm}
		\textbf{SCENARIO:}
		\begin{itemize}
			\item Uber driverless vehicle in Tempe, Arizona
			\item Fatally injured a pedestrian
			\item Multiple failures
		\end{itemize}
		
		\vspace{0.1cm}
		\textbf{FAILURES:}
		\begin{itemize}
			\item \textcolor{myblue}{System Failure:}
			\begin{itemize}
				\item Vehicle's systems failed to recognize pedestrian in time
			\end{itemize}
			\item \textcolor{myblue}{Deactivated Safety:}
			\begin{itemize}
				\item Some safety mechanisms were deactivated
			\end{itemize}
			\item \textcolor{myblue}{Human Factor:}
			\begin{itemize}
				\item Safety driver distracted
				\item Did not react quickly enough
				\item Result of automation bias
			\end{itemize}
		\end{itemize}
		
		\vspace{0.1cm}
		\textbf{LESSONS:}
		\begin{itemize}
			\item Risk of over-dependence on AI
			\item Need for robust testing and validation
			\item Comprehensive testing protocols
		\end{itemize}
		
		\textcolor{myred}{\textbf{EXAM TRAP:}} Uber crash = \textcolor{myblue}{risks of over-reliance on AI}!
	\end{frame}
	
	\begin{frame}{T61: Automation Bias}
		\fontsize{6.5}{7.5}\selectfont
		\textbf{TOPIC: Automation Bias}
		
		\vspace{0.1cm}
		\textbf{DEFINITION:}
		\begin{itemize}
			\item Tendency of humans to over-rely on automated systems
			\item Leading to complacency and reduced vigilance
		\end{itemize}
		
		\vspace{0.1cm}
		\textbf{CONSEQUENCES:}
		\begin{itemize}
			\item Errors of oversight
			\item Especially problematic in critical situations
			\item Skill atrophy
		\end{itemize}
		
		\vspace{0.1cm}
		\textbf{SKILL ATROPHY:}
		\begin{itemize}
			\item Prolonged reliance on automation
			\item Decay in essential skills
			\item Operators become passive observers
			\item From active decision makers to passive
		\end{itemize}
		
		\textcolor{myred}{\textbf{EXAM TRAP:}} Automation bias = \textcolor{myblue}{over-reliance on automated systems}!
	\end{frame}
	
	\begin{frame}{T62: Automation Bias Mitigation}
		\fontsize{6.5}{7.5}\selectfont
		\textbf{TOPIC: Mitigating Automation Bias}
		
		\vspace{0.1cm}
		\textbf{MITIGATION STRATEGIES:}
		\begin{enumerate}
			\item \textcolor{myblue}{Inform Users:}
			\begin{itemize}
				\item Risks of automation bias
				\item Limitations of AI
			\end{itemize}
			\item \textcolor{myblue}{Design Mechanisms:}
			\begin{itemize}
				\item Require periodic human input
				\item Require verification
				\item Keep operators engaged and alert
			\end{itemize}
		\end{enumerate}
		
		\vspace{0.1cm}
		\textbf{EXAMPLE:}
		\begin{itemize}
			\item Eye tracking technology
			\item Combined with alerts
			\item Ensure meaningful human control
			\item Uber could have used this
		\end{itemize}
		
		\textcolor{myred}{\textbf{EXAM TRAP:}} Mitigation = \textcolor{myblue}{user education + periodic human verification}!
	\end{frame}
	
	\begin{frame}{T63: Mental Health AI Risks}
		\fontsize{6.5}{7.5}\selectfont
		\textbf{TOPIC: AI in Mental Health Support}
		
		\vspace{0.1cm}
		\textbf{APPLICATIONS:}
		\begin{itemize}
			\item Chatbots
			\item Virtual therapists
			\item Offer accessibility and anonymity
		\end{itemize}
		
		\vspace{0.1cm}
		\textbf{CONCERNS:}
		\begin{itemize}
			\item Adequacy in handling complex emotional issues
			\item AI has no empathy
			\item Empathy crucial for mental health support
			\item Positive effects dwindle when patients know AI-generated
		\end{itemize}
		
		\vspace{0.1cm}
		\textbf{GOLDEN RULE:}
		\begin{itemize}
			\item Just because we can does not mean we should
			\item Reflect on which issues lend themselves to automation
			\item Some should stay in human hands
		\end{itemize}
		
		\textcolor{myred}{\textbf{EXAM TRAP:}} AI lacks \textcolor{myblue}{genuine empathy} - critical for mental health!
	\end{frame}
	
	\begin{frame}{T64: Golden Rules for AI Development}
		\fontsize{6.5}{7.5}\selectfont
		\textbf{TOPIC: Golden Rules for AI Development}
		
		\vspace{0.1cm}
		\textbf{RULE 1:}
		\begin{itemize}
			\item "Just because we can does not mean we should"
			\item Reflect carefully on which issues lend themselves to automation
			\item Some issues should stay in human hands
			\item Until we have found ways to overcome inadequacies
		\end{itemize}
		
		\vspace{0.1cm}
		\textbf{RULE 2:}
		\begin{itemize}
			\item Involve diverse stakeholders
			\item Domain experts, wider public, ethicists, policy makers
			\item For AI applications in critical situations
			\item Ensures technology is ethical, safe, appropriate
		\end{itemize}
		
		\vspace{0.1cm}
		\textbf{RULE 3:}
		\begin{itemize}
			\item Developers and users must be transparent and accountable
			\item Clear lines of responsibility
			\item In case of failures or adverse outcomes
		\end{itemize}
		
		\textcolor{myred}{\textbf{EXAM TRAP:}} Golden rules = \textcolor{myblue}{reflection, inclusion, transparency}!
	\end{frame}
	
	% ============================================================
	% SECTION 7: REPUTATIONAL RISKS
	% ============================================================
	
	\begin{frame}{T65: Reputational Risks Overview}
		\fontsize{6.5}{7.5}\selectfont
		\textbf{TOPIC: Reputational Risks from AI}
		
		\vspace{0.1cm}
		\textbf{DEFINITION:}
		\begin{itemize}
			\item Companies expose themselves to novel risks
			\item When use or output of algorithms goes against prevalent norms and values
			\item Reputational damage
		\end{itemize}
		
		\vspace{0.1cm}
		\textbf{THREE MAIN CAUSES:}
		\begin{enumerate}
			\item \textcolor{myblue}{Privacy Breaches:}
			\begin{itemize}
				\item Cambridge Analytica example
				\item Mishandling user data
			\end{itemize}
			\item \textcolor{myblue}{Algorithmic Bias:}
			\begin{itemize}
				\item Dutch child benefits scandal
				\item Unfair algorithms
			\end{itemize}
			\item \textcolor{myblue}{Lack of Explainability:}
			\begin{itemize}
				\item Customers denied without reasons
				\item No way to find out why
			\end{itemize}
		\end{enumerate}
		
		\textcolor{myred}{\textbf{EXAM TRAP:}} Three main causes = \textcolor{myblue}{privacy, bias, explainability}!
	\end{frame}
	
	\begin{frame}{T66: Dutch Child Benefits Scandal}
		\fontsize{6.5}{7.5}\selectfont
		\textbf{TOPIC: Dutch Child Benefits Scandal}
		
		\vspace{0.1cm}
		\textbf{SCENARIO:}
		\begin{itemize}
			\item Fraud-detection algorithm deployed by Dutch government
			\item Falsely denied and reclaimed child benefits
			\item Thousands of Dutch citizens affected
			\item Dire consequences for families
		\end{itemize}
		
		\vspace{0.1cm}
		\textbf{OUTCOME:}
		\begin{itemize}
			\item Dutch prime minister resigned
			\item Major political scandal
			\item Reputational damage to government
		\end{itemize}
		
		\vspace{0.1cm}
		\textbf{LESSONS:}
		\begin{itemize}
			\item Algorithmic bias can have severe consequences
			\item Need for oversight and accountability
			\item Importance of fairness in government algorithms
		\end{itemize}
		
		\textcolor{myred}{\textbf{EXAM TRAP:}} Dutch scandal = \textcolor{myblue}{algorithmic bias with severe consequences}!
	\end{frame}
	
	\begin{frame}{T67: Types of AI Criticism}
		\fontsize{6.5}{7.5}\selectfont
		\textbf{TOPIC: Types of AI-Related Criticism}
		
		\vspace{0.1cm}
		\textbf{TWO CATEGORIES:}
		
		\begin{enumerate}
			\item \textcolor{myblue}{Competence Criticism:}
			\begin{itemize}
				\item Perception that company lacks expertise
				\item Insufficient capacities for safe/fair algorithms
				\item Technical failings, unfairness, lack of explainability
				\item Data not stored safely
			\end{itemize}
			
			\item \textcolor{myblue}{Value Alignment Criticism:}
			\begin{itemize}
				\item Company fully aware of risks but proceeded
				\item Facebook emotional contagion experiment
				\item Violation of user trust and expectations
				\item Lack of due diligence
			\end{itemize}
		\end{enumerate}
		
		\vspace{0.1cm}
		\textbf{OVERLAP:}
		\begin{itemize}
			\item Distinction not always clear-cut
			\item Data breach can be incompetence or misaligned values
			\item Perception may be both
		\end{itemize}
		
		\textcolor{myred}{\textbf{EXAM TRAP:}} Competence = \textcolor{myblue}{lack of expertise}; Value alignment = \textcolor{myblue}{misaligned values}!
	\end{frame}
	
	\begin{frame}{T68: Reputational Risk Mitigation}
		\fontsize{6.5}{7.5}\selectfont
		\textbf{TOPIC: Mitigating Reputational Risks}
		
		\vspace{0.1cm}
		\textbf{STRATEGIES:}
		\begin{enumerate}
			\item \textcolor{myblue}{Continuous Monitoring and Evaluation:}
			\begin{itemize}
				\item Identify risks proactively
				\item Organizational governance mechanisms
				\item Regular monitoring and evaluation
			\end{itemize}
			
			\item \textcolor{myblue}{Ensuring Privacy and Security:}
			\begin{itemize}
				\item Comply with GDPR
				\item Secure data storage, handling, processing
				\item Transparent data usage policies
			\end{itemize}
			
			\item \textcolor{myblue}{Ensuring Algorithmic Fairness:}
			\begin{itemize}
				\item Train on diverse datasets
				\item Regular audits for discriminatory patterns
			\end{itemize}
			
			\item \textcolor{myblue}{Promoting Transparency and Explainability:}
			\begin{itemize}
				\item XAI techniques
				\item Customers aware of algorithm use
				\item Builds trust and accountability
			\end{itemize}
		\end{enumerate}
		
		\textcolor{myred}{\textbf{EXAM TRAP:}} Mitigation = \textcolor{myblue}{monitoring, privacy, fairness, transparency}!
	\end{frame}
	
	\begin{frame}{T69: Demonstrating Will to Change}
		\fontsize{6.5}{7.5}\selectfont
		\textbf{TOPIC: Demonstrating Will to Change After Incidents}
		
		\vspace{0.1cm}
		\textbf{DEFINITION:}
		\begin{itemize}
			\item Once incident has happened, limit reputational damage
			\item Be transparent about what went wrong and why
		\end{itemize}
		
		\vspace{0.1cm}
		\textbf{CASES OF INCOMPETENCE:}
		\begin{itemize}
			\item Demonstrate working to fix issue
			\item Show corrective actions
			\item Communicate progress
		\end{itemize}
		
		\vspace{0.1cm}
		\textbf{CASES OF VALUE MISALIGNMENT:}
		\begin{itemize}
			\item Problem rooted in governance structures
			\item Or company culture
			\item Reassure stakeholders
			\item Fix underlying problems via organizational means
			\item Transparency key to rebuilding trust
		\end{itemize}
		
		\textcolor{myred}{\textbf{EXAM TRAP:}} Transparency is \textcolor{myblue}{key to rebuilding trust}!
	\end{frame}
	
	\begin{frame}{T70: Ethical AI Frameworks}
		\fontsize{6.5}{7.5}\selectfont
		\textbf{TOPIC: Ethical AI Frameworks}
		
		\vspace{0.1cm}
		\textbf{DEFINITION:}
		\begin{itemize}
			\item Ethical guidelines to prevent incidents
			\item Provide guidance and overview of risks
		\end{itemize}
		
		\vspace{0.1cm}
		\textbf{BENEFITS:}
		\begin{itemize}
			\item Prevent incidents
			\item Provide guidance
			\item Overview of risks
		\end{itemize}
		
		\vspace{0.1cm}
		\textbf{CHALLENGES:}
		\begin{itemize}
			\item Many frameworks too vague
			\item Not helpful in cases of ambiguity
			\item Developing in collaboration with experts helps
			\item Ensure they provide needed guidance
		\end{itemize}
		
		\textcolor{myred}{\textbf{EXAM TRAP:}} Ethical frameworks must be \textcolor{myblue}{specific enough to be helpful}!
	\end{frame}
	
	\begin{frame}{T71: Stakeholder Engagement}
		\fontsize{6.5}{7.5}\selectfont
		\textbf{TOPIC: Stakeholder Engagement and Communication}
		
		\vspace{0.1cm}
		\textbf{DEFINITION:}
		\begin{itemize}
			\item Regularly engage with stakeholders
			\item Customers, employees, regulators
			\item About AI initiatives and their impacts
		\end{itemize}
		
		\vspace{0.1cm}
		\textbf{BENEFITS:}
		\begin{itemize}
			\item Clear communication prevents misunderstandings
			\item Builds trust
			\item Early identification of concerns
			\item Collaborative problem-solving
		\end{itemize}
		
		\vspace{0.1cm}
		\textbf{APPROACH:}
		\begin{itemize}
			\item Regular updates
			\item Transparent about challenges
			\item Responsive to feedback
			\item Inclusive of diverse perspectives
		\end{itemize}
		
		\textcolor{myred}{\textbf{EXAM TRAP:}} Stakeholder engagement \textcolor{myblue}{builds trust and prevents misunderstandings}!
	\end{frame}
	
	\begin{frame}{T72: Crisis Management Plan}
		\fontsize{6.5}{7.5}\selectfont
		\textbf{TOPIC: AI-Related Crisis Management Plan}
		
		\vspace{0.1cm}
		\textbf{DEFINITION:}
		\begin{itemize}
			\item Perform risk analysis
			\item Inform development of crisis-management plan
			\item Tailored to AI-related incidents
		\end{itemize}
		
		\vspace{0.1cm}
		\textbf{COMPONENTS:}
		\begin{itemize}
			\item Communication strategies
			\item Potential steps to rectify issue
			\item Clear lines of responsibility
			\item Escalation procedures
		\end{itemize}
		
		\vspace{0.1cm}
		\textbf{BENEFITS:}
		\begin{itemize}
			\item Faster response to incidents
			\item Minimize reputational damage
			\item Demonstrate competence and responsibility
			\item Rebuild trust more effectively
		\end{itemize}
		
		\textcolor{myred}{\textbf{EXAM TRAP:}} Crisis management plan = \textcolor{myblue}{proactive preparation}!
	\end{frame}
	
	% ============================================================
	% SECTION 8: EXISTENTIAL AND GLOBAL RISKS
	% ============================================================
	
	\begin{frame}{T73: Existential Risks}
		\fontsize{6.5}{7.5}\selectfont
		\textbf{TOPIC: Existential Risks from AI}
		
		\vspace{0.1cm}
		\textbf{DEFINITION:}
		\begin{itemize}
			\item Scenarios where advanced AI poses severe or catastrophic threats to human existence
			\item Central theme: superintelligence
		\end{itemize}
		
		\vspace{0.1cm}
		\textbf{SUPERINTELLIGENCE:}
		\begin{itemize}
			\item AI system operating beyond human-level intelligence
			\item Could act in ways that threaten human society or existence
			\item Concern about well-intentioned AI with misaligned objectives
		\end{itemize}
		
		\vspace{0.1cm}
		\textbf{KEY CONCERNS:}
		\begin{itemize}
			\item \textcolor{myblue}{Rapid Pace:} AI development follows rapid pace
			\item \textcolor{myblue}{AI Alignment:} Ensuring means and goals aligned with human values
			\item \textcolor{myblue}{AI Arms Race:} Different stakeholders creating powerful AI without safety measures
		\end{itemize}
		
		\textcolor{myred}{\textbf{EXAM TRAP:}} Existential risks = \textcolor{myblue}{threats to human existence}!
	\end{frame}
	
	\begin{frame}{T74: AI Alignment}
		\fontsize{6.5}{7.5}\selectfont
		\textbf{TOPIC: AI Alignment}
		
		\vspace{0.1cm}
		\textbf{DEFINITION:}
		\begin{itemize}
			\item Ensuring AI systems' means and goals align with human values
			\item Central concern for existential risk
		\end{itemize}
		
		\vspace{0.1cm}
		\textbf{EXAMPLE:}
		\begin{itemize}
			\item Paperclip maximizer scenario
			\item Superintelligent AI programmed to make paperclips
			\item Wipes out humanity because helps make more paperclips
			\item Well-intentioned AI with misaligned objective
		\end{itemize}
		
		\vspace{0.1cm}
		\textbf{CHALLENGES:}
		\begin{itemize}
			\item Difficulty specifying human values
			\item Unintended consequences
			\item Value lock-in
			\item Ensuring AI understands what we actually want
		\end{itemize}
		
		\textcolor{myred}{\textbf{EXAM TRAP:}} AI alignment = \textcolor{myblue}{matching AI to human values}!
	\end{frame}
	
	\begin{frame}{T75: AI Arms Race}
		\fontsize{6.5}{7.5}\selectfont
		\textbf{TOPIC: AI Arms Race}
		
		\vspace{0.1cm}
		\textbf{DEFINITION:}
		\begin{itemize}
			\item Risks from different stakeholders creating powerful AI
			\item Without adequate safety measures
			\item Could lead to catastrophic global conflicts
		\end{itemize}
		
		\vspace{0.1cm}
		\textbf{CHARACTERISTICS:}
		\begin{itemize}
			\item Competition between nations/companies
			\item Pressure to develop quickly
			\item Safety measures may be sacrificed
			\item Potential for existential outcomes
		\end{itemize}
		
		\vspace{0.1cm}
		\textbf{COUNTERARGUMENTS:}
		\begin{itemize}
			\item Current AI is "narrow"
			\item Far from superintelligent
			\item Regulatory mechanisms will develop
			\item Historical parallels to hype
		\end{itemize}
		
		\textcolor{myred}{\textbf{EXAM TRAP:}} AI arms race = \textcolor{myblue}{unsafe AI development competition}!
	\end{frame}
	
	\begin{frame}{T76: Economic Risks from AI}
		\fontsize{6.5}{7.5}\selectfont
		\textbf{TOPIC: Economic Risks from AI}
		
		\vspace{0.1cm}
		\textbf{JOB DISPLACEMENT:}
		\begin{itemize}
			\item Initial study (2017): 47\% of US employment at risk
			\item Subsequent studies: lower numbers
			\item ChatGPT reignited discussion
			\item Accenture: 40\% of working hours could be affected
			\item WEF (2023): 83 million jobs destroyed, 69 million created
			\item Goldman Sachs (2023): 2/3 jobs exposed, 1/4 automated
		\end{itemize}
		
		\vspace{0.1cm}
		\textbf{CHARACTERISTICS:}
		\begin{itemize}
			\item High-skill jobs also at stake (accounting)
			\item Gap between AI-skilled and unskilled
			\item Widening socio-economic inequalities
		\end{itemize}
		
		\vspace{0.1cm}
		\textbf{CAUTION:}
		\begin{itemize}
			\item Predictions should be taken with care
			\item Difficult to predict pace of AI advancement
			\item Human supervision requirements unclear
		\end{itemize}
		
		\textcolor{myred}{\textbf{EXAM TRAP:}} Economic risks = \textcolor{myblue}{job displacement and inequality}!
	\end{frame}
	
	\begin{frame}{T77: Global Inequality}
		\fontsize{6.5}{7.5}\selectfont
		\textbf{TOPIC: Global Inequality from AI}
		
		\vspace{0.1cm}
		\textbf{DEFINITION:}
		\begin{itemize}
			\item Most large AI firms based in US
			\item Technological adoption varies widely
			\item Risk of exacerbating global inequalities
		\end{itemize}
		
		\vspace{0.1cm}
		\textbf{CONSEQUENCES:}
		\begin{itemize}
			\item Nations at center of AI development leap ahead
			\item In economic growth and political influence
			\item Less technologically developed countries left behind
		\end{itemize}
		
		\vspace{0.1cm}
		\textbf{IMPLICATIONS:}
		\begin{itemize}
			\item Ethical considerations
			\item Risks of economic and geopolitical tensions
			\item Need for international cooperation
			\item Equitable distribution of AI benefits
		\end{itemize}
		
		\textcolor{myred}{\textbf{EXAM TRAP:}} AI could \textcolor{myblue}{widen global inequality}!
	\end{frame}
	
	\begin{frame}{T78: Misinformation Campaigns}
		\fontsize{6.5}{7.5}\selectfont
		\textbf{TOPIC: AI-Enabled Misinformation Campaigns}
		
		\vspace{0.1cm}
		\textbf{DEFINITION:}
		\begin{itemize}
			\item AI enables sophisticated misinformation campaigns
			\item Bots creating misleading social media posts
			\item AI-generated imagery and videos
		\end{itemize}
		
		\vspace{0.1cm}
		\textbf{RISKS:}
		\begin{itemize}
			\item Mislead and influence public opinion
			\item Threat to political institutions
			\item Threat to democratic elections
			\item Loss of trust in information
		\end{itemize}
		
		\vspace{0.1cm}
		\textbf{MANAGEMENT:}
		\begin{itemize}
			\item Strict governance measures
			\item Ensure information authenticity and reliability
			\item AI detection tools
			\item Expose AI-generated content
		\end{itemize}
		
		\textcolor{myred}{\textbf{EXAM TRAP:}} Misinformation = \textcolor{myblue}{threat to democratic institutions}!
	\end{frame}
	
	\begin{frame}{T79: Privacy and Surveillance}
		\fontsize{6.5}{7.5}\selectfont
		\textbf{TOPIC: Privacy and Surveillance Risks}
		
		\vspace{0.1cm}
		\textbf{DEFINITION:}
		\begin{itemize}
			\item AI models require large amounts of data for training
			\item Strong motivation for companies to collect data
			\item Risks to privacy exploded with data economy
		\end{itemize}
		
		\vspace{0.1cm}
		\textbf{RISKS:}
		\begin{itemize}
			\item Mass surveillance
			\item Data extraction without consent
			\item Profiling without knowledge
			\item Manipulation through targeting
		\end{itemize}
		
		\vspace{0.1cm}
		\textbf{REGULATORY RESPONSES:}
		\begin{itemize}
			\item GDPR
			\item Data protection laws
			\item Transparency requirements
			\item Consent requirements
		\end{itemize}
		
		\textcolor{myred}{\textbf{EXAM TRAP:}} Privacy risks \textcolor{myblue}{exploded with data economy}!
	\end{frame}
	
	% ============================================================
	% SECTION 9: GENAI RISKS
	% ============================================================
	
	\begin{frame}{T80: GenAI Overview}
		\fontsize{6.5}{7.5}\selectfont
		\textbf{TOPIC: Generative AI Overview}
		
		\vspace{0.1cm}
		\textbf{DEFINITION:}
		\begin{itemize}
			\item Creates new content: text, audio, images, videos
			\item Used across industries
			\item High-volume, text-based processes
		\end{itemize}
		
		\vspace{0.1cm}
		\textbf{ADOPTION STATISTICS:}
		\begin{itemize}
			\item 71\% of organizations regularly use GenAI (McKinsey 2025)
			\item Up from 33\% in 2023
			\item Banking and financial services: 58\% embrace GenAI
			\item Professional services: 95\% use at least once a month (BCG)
		\end{itemize}
		
		\vspace{0.1cm}
		\textbf{APPLICATIONS:}
		\begin{itemize}
			\item Customer communication
			\item Marketing
			\item Internal documentation
			\item Code generation
			\item Decision-support systems
		\end{itemize}
		
		\textcolor{myred}{\textbf{EXAM TRAP:}} GenAI adoption = \textcolor{myblue}{71\% of organizations} (2025 McKinsey)!
	\end{frame}
	
	\begin{frame}{T81: GenAI Risk Profile}
		\fontsize{6.5}{7.5}\selectfont
		\textbf{TOPIC: GenAI Risk Profile - Five Domains}
		
		\vspace{0.1cm}
		\textbf{FIVE RISK DOMAINS:}
		\begin{enumerate}
			\item \textcolor{myblue}{Functional:} Performance and reliability
			\item \textcolor{myblue}{Data-Related:} Training, fine-tuning, access to sensitive data
			\item \textcolor{myblue}{Organizational:} Adoption, embedding, management
			\item \textcolor{myblue}{Strategic/Systemic:} Long-term positioning, dependencies
			\item \textcolor{myblue}{Societal:} Democratic institutions, public discourse
		\end{enumerate}
		
		\vspace{0.1cm}
		\textbf{KEY CHARACTERISTICS:}
		\begin{itemize}
			\item Risks not only technical
			\item Organizational and societal
			\item Model behavior can change through updates
			\item How model is used matters
		\end{itemize}
		
		\textcolor{myred}{\textbf{EXAM TRAP:}} GenAI risks span \textcolor{myblue}{five domains}!
	\end{frame}
	
	\begin{frame}{T82: Hallucinations}
		\fontsize{6.5}{7.5}\selectfont
		\textbf{TOPIC: Hallucinations in GenAI}
		
		\vspace{0.1cm}
		\textbf{DEFINITION:}
		\begin{itemize}
			\item Generating factually incorrect, misleading, or fabricated outputs
			\item Appear credible
			\item Difficult for non-expert users to detect
		\end{itemize}
		
		\vspace{0.1cm}
		\textbf{CONSEQUENCES:}
		\begin{itemize}
			\item Reputational damages
			\item Legal liability
			\item Misinformation
		\end{itemize}
		
		\vspace{0.1cm}
		\textbf{EXAMPLE:}
		\begin{itemize}
			\item Air Canada chatbot offered discount
			\item Discount shouldn't have been available
			\item Company held responsible
			\item Awarded damages and tribunal fees
		\end{itemize}
		
		\textcolor{myred}{\textbf{EXAM TRAP:}} Hallucination = \textcolor{myblue}{incorrect but credible outputs}!
	\end{frame}
	
	\begin{frame}{T83: Unpredictability}
		\fontsize{6.5}{7.5}\selectfont
		\textbf{TOPIC: Unpredictability in GenAI}
		
		\vspace{0.1cm}
		\textbf{DEFINITION:}
		\begin{itemize}
			\item GenAI systems unpredictable in output
			\item Capabilities develop unpredictably
			\item Undermines standardization
		\end{itemize}
		
		\vspace{0.1cm}
		\textbf{CHALLENGES:}
		\begin{itemize}
			\item Standardize outputs across users and time
			\item Auditability
			\item Consistent treatment
			\item Compliance
		\end{itemize}
		
		\vspace{0.1cm}
		\textbf{APPLICATIONS AFFECTED:}
		\begin{itemize}
			\item Financial advice
			\item Claims handling
			\item Legal contexts
			\item Anywhere consistency required
		\end{itemize}
		
		\textcolor{myred}{\textbf{EXAM TRAP:}} Unpredictability undermines \textcolor{myblue}{standardization and auditability}!
	\end{frame}
	
	\begin{frame}{T84: Adversarial Inputs}
		\fontsize{6.5}{7.5}\selectfont
		\textbf{TOPIC: Adversarial Inputs in GenAI}
		
		\vspace{0.1cm}
		\textbf{DEFINITION:}
		\begin{itemize}
			\item Techniques to cause unintended behavior
			\item Exploit vulnerabilities in GenAI systems
		\end{itemize}
		
		\vspace{0.1cm}
		\textbf{TYPES:}
		\begin{itemize}
			\item \textcolor{myblue}{Prompt Injection:}
			\begin{itemize}
				\item Hidden instructions within user inputs
				\item Cause unintended behavior
			\end{itemize}
			\item \textcolor{myblue}{Jailbreaks:}
			\begin{itemize}
				\item Bypassing safety mechanisms
				\item Crafted prompts
			\end{itemize}
			\item \textcolor{myblue}{Other Manipulation:}
			\begin{itemize}
				\item Various techniques
				\item Cause unsafe behavior
			\end{itemize}
		\end{itemize}
		
		\textcolor{myred}{\textbf{EXAM TRAP:}} Adversarial inputs = \textcolor{myblue}{hidden instructions to cause unintended behavior}!
	\end{frame}
	
	\begin{frame}{T85: Data Privacy Violations}
		\fontsize{6.5}{7.5}\selectfont
		\textbf{TOPIC: Data Privacy Violations in GenAI}
		
		\vspace{0.1cm}
		\textbf{DEFINITION:}
		\begin{itemize}
			\item GenAI relies on user input
			\item Prompts may contain personally identifiable information
			\item Data may be inadvertently stored, shared, or used for training
		\end{itemize}
		
		\vspace{0.1cm}
		\textbf{RISKS:}
		\begin{itemize}
			\item Without user's knowledge
			\item Especially with third-party providers
			\item If no transparency on how inputs handled
		\end{itemize}
		
		\vspace{0.1cm}
		\textbf{EXAMPLES:}
		\begin{itemize}
			\item Researchers extracted emails, phone numbers from LLM
			\item FTC investigation into OpenAI
			\item Data privacy issues
			\item Violation of consumer protection laws
		\end{itemize}
		
		\textcolor{myred}{\textbf{EXAM TRAP:}} GenAI prompts may \textcolor{myblue}{inadvertently expose PII}!
	\end{frame}
	
	\begin{frame}{T86: Copyright Issues}
		\fontsize{6.5}{7.5}\selectfont
		\textbf{TOPIC: Copyright Issues in GenAI}
		
		\vspace{0.1cm}
		\textbf{DEFINITION:}
		\begin{itemize}
			\item GenAI models trained on vast web-scraped content
			\item Includes copyrighted books, images, articles, software
			\item Creates debate on copyright infringements
		\end{itemize}
		
		\vspace{0.1cm}
		\textbf{LEGAL PROCEEDINGS:}
		\begin{itemize}
			\item Getty Images vs. Stability AI
			\item Infringing UK intellectual property rights
			\item Main copyright claim dropped
			\item Other claims continue (secondary infringement, trademarks)
		\end{itemize}
		
		\vspace{0.1cm}
		\textbf{IMPLICATIONS:}
		\begin{itemize}
			\item Using generated content may carry copyright risks
			\item Organizations must consider
			\item Public-facing content especially
		\end{itemize}
		
		\textcolor{myred}{\textbf{EXAM TRAP:}} GenAI training may \textcolor{myblue}{infringe copyright}!
	\end{frame}
	
	\begin{frame}{T87: Responsibility Gaps}
		\fontsize{6.5}{7.5}\selectfont
		\textbf{TOPIC: Responsibility Gaps in GenAI}
		
		\vspace{0.1cm}
		\textbf{DEFINITION:}
		\begin{itemize}
			\item GenAI systems fall into gap between existing lines of accountability
			\item Adopted by everyone, including non-technical teams
			\item HR team using LLMs to screen complaints
		\end{itemize}
		
		\vspace{0.1cm}
		\textbf{PROBLEM:}
		\begin{itemize}
			\item Biased or unfair outputs possible
			\item Need clear lines of responsibility
			\item Quality assurance
			\item Oversight and risk mitigation
		\end{itemize}
		
		\vspace{0.1cm}
		\textbf{SOLUTION:}
		\begin{itemize}
			\item Establish clear lines of responsibility
			\item For GenAI generated outputs
			\item Both quality assurance and risk mitigation
		\end{itemize}
		
		\textcolor{myred}{\textbf{EXAM TRAP:}} GenAI adoption creates \textcolor{myblue}{responsibility gaps}!
	\end{frame}
	
	\begin{frame}{T88: Misaligned Incentives}
		\fontsize{6.5}{7.5}\selectfont
		\textbf{TOPIC: Misaligned Incentives in GenAI}
		
		\vspace{0.1cm}
		\textbf{DEFINITION:}
		\begin{itemize}
			\item Pressure to innovate and automate quickly
			\item Drive to cut costs outweighs cautionary approaches
			\item Considerations about reliability, safety, long-term risk
		\end{itemize}
		
		\vspace{0.1cm}
		\textbf{EXAMPLES:}
		\begin{itemize}
			\item Anthropic lobbied to change safety bill
			\item Focus on post-hoc accountability
			\item Teams rewarded for cost savings
			\item Not for flagging risks
		\end{itemize}
		
		\vspace{0.1cm}
		\textbf{CONSEQUENCE:}
		\begin{itemize}
			\item Incentives favor rapid deployment
			\item Likely to be risky
			\item Over safe deployment
		\end{itemize}
		
		\textcolor{myred}{\textbf{EXAM TRAP:}} Innovation pressure can \textcolor{myblue}{outweigh safety considerations}!
	\end{frame}
	
	\begin{frame}{T89: Lack of Training}
		\fontsize{6.5}{7.5}\selectfont
		\textbf{TOPIC: Lack of Internal Guidance or Training}
		
		\vspace{0.1cm}
		\textbf{DEFINITION:}
		\begin{itemize}
			\item Lack of formal policy
			\item Lack of training of employees
			\item Use GenAI systems without guidance
		\end{itemize}
		
		\vspace{0.1cm}
		\textbf{RISKS:}
		\begin{itemize}
			\item Employees integrate into workflow
			\item Without understanding acceptable use
			\item Data protection obligations
			\item Reputational risks
			\item Lines of accountability
		\end{itemize}
		
		\vspace{0.1cm}
		\textbf{EXAMPLE:}
		\begin{itemize}
			\item Cyberhaven: 4.2\% of workers pasted sensitive data into ChatGPT
			\item JPMorgan, Amazon, Verizon, Accenture
			\item Banned ChatGPT for employees
			\item Compliance concerns
		\end{itemize}
		
		\textcolor{myred}{\textbf{EXAM TRAP:}} Lack of training leads to \textcolor{myblue}{unintentional data exposure}!
	\end{frame}
	
	\begin{frame}{T90: Vendor Lock-In}
		\fontsize{6.5}{7.5}\selectfont
		\textbf{TOPIC: Vendor Lock-In in GenAI}
		
		\vspace{0.1cm}
		\textbf{DEFINITION:}
		\begin{itemize}
			\item Reliance on external providers
			\item Small set of dominant suppliers
			\item Creates dependencies and power asymmetries
		\end{itemize}
		
		\vspace{0.1cm}
		\textbf{RISKS:}
		\begin{itemize}
			\item Limit auditability, transparency, leverage
			\item Providers change pricing, models, usage
			\item Without warning
			\item Minimal control over model behavior
		\end{itemize}
		
		\vspace{0.1cm}
		\textbf{CONSEQUENCE:}
		\begin{itemize}
			\item Short-term adoption decisions
			\item Shift into long-term structural lock-ins
			\item Reduced flexibility
		\end{itemize}
		
		\textcolor{myred}{\textbf{EXAM TRAP:}} Vendor lock-in = \textcolor{myblue}{dependence limiting control}!
	\end{frame}
	
	\begin{frame}{T91: Regulatory Unpredictability}
		\fontsize{6.5}{7.5}\selectfont
		\textbf{TOPIC: Regulatory Unpredictability in GenAI}
		
		\vspace{0.1cm}
		\textbf{DEFINITION:}
		\begin{itemize}
			\item Regulation evolving quickly but unevenly
			\item EU AI Act: Ex ante obligations
			\item US: Patchwork of state laws
		\end{itemize}
		
		\vspace{0.1cm}
		\textbf{CHALLENGES:}
		\begin{itemize}
			\item Organizations moving across jurisdictions
			\item Even low-risk use cases subject to complex requirements
			\item Future requirements uncertain
		\end{itemize}
		
		\vspace{0.1cm}
		\textbf{EXAMPLE:}
		\begin{itemize}
			\item Apple delayed AI features in Europe
			\item Citing regulatory uncertainty
			\item Compliance concerns
			\item Cautious approach
		\end{itemize}
		
		\textcolor{myred}{\textbf{EXAM TRAP:}} Regulatory unpredictability = \textcolor{myblue}{compliance challenges across jurisdictions}!
	\end{frame}
	
	\begin{frame}{T92: Public Trust and Reputation}
		\fontsize{6.5}{7.5}\selectfont
		\textbf{TOPIC: Public Trust and Reputation Risks in GenAI}
		
		\vspace{0.1cm}
		\textbf{DEFINITION:}
		\begin{itemize}
			\item GenAI systems public-facing
			\item Draft emails, respond to queries, generate content
			\item Biased content, hallucinations, failure to disclose
		\end{itemize}
		
		\vspace{0.1cm}
		\textbf{RISKS:}
		\begin{itemize}
			\item Public backlash
			\item Even if compliant with regulations
			\item Reputation damage
			\item Regulatory scrutiny
		\end{itemize}
		
		\vspace{0.1cm}
		\textbf{EXAMPLE:}
		\begin{itemize}
			\item Grok chatbot (X/Twitter)
			\item Started using Nazi rhetoric
			\item Insulted public leaders
			\item Global outrage
			\item Turkey banned Grok
		\end{itemize}
		
		\textcolor{myred}{\textbf{EXAM TRAP:}} Public trust = \textcolor{myblue}{critical for GenAI deployment}!
	\end{frame}
	
	\begin{frame}{T93: Disinformation at Scale}
		\fontsize{6.5}{7.5}\selectfont
		\textbf{TOPIC: Disinformation and Synthetic Content at Scale}
		
		\vspace{0.1cm}
		\textbf{DEFINITION:}
		\begin{itemize}
			\item GenAI lowers cost of creating misleading content
			\item Fabricated news, health misinformation
			\item Political content
		\end{itemize}
		
		\vspace{0.1cm}
		\textbf{RISKS:}
		\begin{itemize}
			\item Influence campaigns
			\item Voter suppression
			\item Damage democratic institutions
			\item Distrust of legitimate sources
		\end{itemize}
		
		\vspace{0.1cm}
		\textbf{REGULATORY RESPONSE:}
		\begin{itemize}
			\item EU Code of Practice on Disinformation
			\item Flags these risks
			\item Need for detection tools
		\end{itemize}
		
		\textcolor{myred}{\textbf{EXAM TRAP:}} Synthetic content can \textcolor{myblue}{damage trust and institutions}!
	\end{frame}
	
	\begin{frame}{T94: Deepfakes}
		\fontsize{6.5}{7.5}\selectfont
		\textbf{TOPIC: Deepfakes and Impersonation}
		
		\vspace{0.1cm}
		\textbf{DEFINITION:}
		\begin{itemize}
			\item Audio or video deepfakes
			\item Convincingly mimic public figures, journalists, private individuals
			\item Individual and societal harm
		\end{itemize}
		
		\vspace{0.1cm}
		\textbf{RISKS:}
		\begin{itemize}
			\item Fraud
			\item Defamation
			\item Political manipulation
			\item Even low-effort impersonations can go viral
		\end{itemize}
		
		\vspace{0.1cm}
		\textbf{EXAMPLE:}
		\begin{itemize}
			\item CEO of UK energy company
			\item Persuaded to transfer €220,000
			\item Voice deepfakes
			\item Criminal use
		\end{itemize}
		
		\textcolor{myred}{\textbf{EXAM TRAP:}} Deepfakes can \textcolor{myblue}{cause fraud, defamation, and manipulation}!
	\end{frame}
	
	\begin{frame}{T95: Manipulative Targeting}
		\fontsize{6.5}{7.5}\selectfont
		\textbf{TOPIC: Manipulative Targeting and Influence}
		
		\vspace{0.1cm}
		\textbf{DEFINITION:}
		\begin{itemize}
			\item Combined with profiling or microtargeting
			\item Especially on social media
			\item GenAI generates persuasive content
		\end{itemize}
		
		\vspace{0.1cm}
		\textbf{CHARACTERISTICS:}
		\begin{itemize}
			\item Tailored to individuals' emotions
			\item Beliefs, personalities
			\item Exacerbates concerns about AI-driven manipulation
		\end{itemize}
		
		\vspace{0.1cm}
		\textbf{RISKS:}
		\begin{itemize}
			\item Exploitation
			\item Undue influence
			\item Erosion of autonomy
			\item Democratic processes affected
		\end{itemize}
		
		\textcolor{myred}{\textbf{EXAM TRAP:}} Manipulative targeting = \textcolor{myblue}{exploitation through personalization}!
	\end{frame}
	
	\begin{frame}{T96: Lack of Provenance}
		\fontsize{6.5}{7.5}\selectfont
		\textbf{TOPIC: Lack of Provenance and Content Traceability}
		
		\vspace{0.1cm}
		\textbf{DEFINITION:}
		\begin{itemize}
			\item Without watermarking or content labels
			\item Synthetic content difficult to identify or regulate
		\end{itemize}
		
		\vspace{0.1cm}
		\textbf{CHALLENGES:}
		\begin{itemize}
			\item Platforms cannot distinguish legitimate from synthetic
			\item Regulators cannot regulate effectively
			\item Users cannot make informed decisions
		\end{itemize}
		
		\vspace{0.1cm}
		\textbf{SOLUTIONS:}
		\begin{itemize}
			\item Watermarking
			\item Content labeling
			\item Provenance tracking
			\item Detection tools
		\end{itemize}
		
		\textcolor{myred}{\textbf{EXAM TRAP:}} Lack of provenance makes \textcolor{myblue}{synthetic content difficult to identify}!
	\end{frame}
	
	\begin{frame}{T97: GenAI Risk Management}
		\fontsize{6.5}{7.5}\selectfont
		\textbf{TOPIC: GenAI Risk Management Strategies}
		
		\vspace{0.1cm}
		\textbf{SHORT-TERM:}
		\begin{itemize}
			\item Existing AI-focused risk management frameworks
			\item Basic policies on GenAI use
			\item Oversight and escalation procedures
			\item Internal review boards
			\item "Responsible AI" champions
		\end{itemize}
		
		\vspace{0.1cm}
		\textbf{MEDIUM-TERM:}
		\begin{itemize}
			\item Align incentives with responsible deployment
			\item Supplement speed and cost metrics
			\item Add risk and reliability metrics
			\item Cross-team collaboration
			\item IT, legal, product teams coordination
		\end{itemize}
		
		\vspace{0.1cm}
		\textbf{KEY PRINCIPLES:}
		\begin{itemize}
			\item Risk identification
			\item Response coordination
			\item Clear accountability
		\end{itemize}
		
		\textcolor{myred}{\textbf{EXAM TRAP:}} Risk management = \textcolor{myblue}{short-term + medium-term strategies}!
	\end{frame}
	
	\begin{frame}{T98: Internal Review Boards}
		\fontsize{6.5}{7.5}\selectfont
		\textbf{TOPIC: Internal Review Boards for AI}
		
		\vspace{0.1cm}
		\textbf{DEFINITION:}
		\begin{itemize}
			\item Create internal review boards
			\item Or designate "responsible AI" champions
			\item Within organization
		\end{itemize}
		
		\vspace{0.1cm}
		\textbf{RESPONSIBILITIES:}
		\begin{itemize}
			\item Raise awareness
			\item Train staff
			\item Guide implementation
			\item GenAI risk management
			\item Review high-risk use cases
		\end{itemize}
		
		\vspace{0.1cm}
		\textbf{BENEFITS:}
		\begin{itemize}
			\item Centralized expertise
			\item Consistent approach
			\item Accountability
			\item Knowledge sharing
		\end{itemize}
		
		\textcolor{myred}{\textbf{EXAM TRAP:}} Review boards \textcolor{myblue}{guide responsible AI implementation}!
	\end{frame}
	
	\begin{frame}{T99: Aligning Incentives}
		\fontsize{6.5}{7.5}\selectfont
		\textbf{TOPIC: Aligning Incentives with Responsible AI}
		
		\vspace{0.1cm}
		\textbf{DEFINITION:}
		\begin{itemize}
			\item Align incentives and reviews
			\item With responsible deployment practices
		\end{itemize}
		
		\vspace{0.1cm}
		\textbf{APPROACH:}
		\begin{itemize}
			\item Supplement speed and cost metrics
			\item Add risk and reliability metrics
			\item Reward responsible behavior
			\item Not just cost savings
		\end{itemize}
		
		\vspace{0.1cm}
		\textbf{CROSS-TEAM COLLABORATION:}
		\begin{itemize}
			\item IT, legal, product teams
			\item Coordination on risk identification
			\item Coordination on responses
			\item Coordination on accountability
		\end{itemize}
		
		\textcolor{myred}{\textbf{EXAM TRAP:}} Align incentives = \textcolor{myblue}{reward responsible AI practices}!
	\end{frame}
	
	\begin{frame}{T100: Module Summary}
		\fontsize{6.5}{7.5}\selectfont
		\textbf{TOPIC: Module 3 Summary - Key Takeaways}
		
		\vspace{0.1cm}
		\textbf{ALGORITHMIC BIAS AND FAIRNESS:}
		\begin{itemize}
			\item Individual vs. group fairness
			\item Three key measures: demographic parity, predictive rate parity, equal opportunity
			\item Fairness impossibility theorem
			\item Multiple sources of unfairness
		\end{itemize}
		
		\vspace{0.1cm}
		\textbf{EXPLAINABILITY:}
		\begin{itemize}
			\item Explainability, interpretability, transparency
			\item XAI techniques: feature importance, Shapley values, LUCID, surrogate models
			\item Black-box problem and opaqueness
		\end{itemize}
		
		\vspace{0.1cm}
		\textbf{AUTONOMY AND SAFETY:}
		\begin{itemize}
			\item Human autonomy and manipulation
			\item Automation bias and skill atrophy
			\item AI in critical systems
		\end{itemize}
		
		\vspace{0.1cm}
		\textbf{EXISTENTIAL AND GLOBAL RISKS:}
		\begin{itemize}
			\item Superintelligence and AI alignment
			\item Economic risks and global inequality
			\item Misinformation and privacy
		\end{itemize}
		
		\vspace{0.1cm}
		\textbf{GENAI RISKS:}
		\begin{itemize}
			\item Five risk domains: functional, data, organizational, strategic, societal
			\item Hallucinations, unpredictability, adversarial inputs
			\item Privacy, copyright, responsibility gaps
			\item Deepfakes, disinformation, manipulation
		\end{itemize}
		
		\textcolor{myred}{\textbf{EXAM SUCCESS:}} Understand all risk categories and their interconnections!
	\end{frame}
	
\end{document}